% =====================================================================
%  LDA-T318 — Lecture 2 — Pretraining
%  Thursday 29 October 2026, 16:15–17:45
%  Build:  latexmk -xelatex main.tex
% =====================================================================
\documentclass[aspectratio=169,10pt]{beamer}

% Path has no parent-directory step: Overleaf resolves \input from the project root and
% will not read above it. Locally, the Makefile puts common/ on TEXINPUTS.
% =====================================================================
%  LDA-T318 — Large Language Models and Generative AI for NLP
%  Shared Beamer preamble for all lecture decks.
%  Aarne Talman, University of Helsinki, 2026.
%
%  Engine:  XeLaTeX or LuaLaTeX preferred (gives the Fira font set).
%           pdfLaTeX also works; Metropolis falls back to Latin Modern.
%  Build:   latexmk -xelatex main.tex      (from the deck's own directory)
%
%  Every deck does:
%      \documentclass[aspectratio=169,10pt]{beamer}
%      \input{../common/preamble.tex}
% =====================================================================

% ---------------------------------------------------------------- theme
\usetheme{metropolis}
\metroset{
  titleformat=smallcaps,
  numbering=fraction,
  progressbar=frametitle,
  block=fill,
  sectionpage=progressbar,
  subsectionpage=none,
}

% ---------------------------------------------------------------- packages
% inputenc/fontenc are required under pdfLaTeX and actively harmful under
% XeLaTeX/LuaLaTeX, where Metropolis loads fontspec to get Fira. Guard them.
\usepackage{iftex}
\ifPDFTeX
  \usepackage[utf8]{inputenc}
  \usepackage[T1]{fontenc}
  \usepackage[sfdefault]{FiraSans}
  \usepackage{FiraMono}
\else
  % Metropolis only gets Fira if fontspec can actually find it. TeX Live ships
  % the OTFs but does NOT register them with fontconfig, so a plain
  % \setsansfont{Fira Sans} silently falls back to Latin Modern and the deck
  % stops looking like Metropolis. Load them by filename instead.
  \usepackage{fontspec}
  \setsansfont{FiraSans}[
    Extension      = .otf,
    UprightFont    = *-Book,
    ItalicFont     = *-BookItalic,
    BoldFont       = *-Medium,
    BoldItalicFont = *-MediumItalic]
  \setmonofont{FiraMono}[
    Extension   = .otf,
    UprightFont = *-Regular,
    BoldFont    = *-Medium,
    Scale       = 0.92]
  \renewcommand{\familydefault}{\sfdefault}
\fi
\usepackage[english]{babel}

\usepackage{amsmath,amssymb}
\usepackage{booktabs}
\usepackage{array}
\usepackage{multicol}
\usepackage{xcolor}
\usepackage{graphicx}
\usepackage{tikz}
\usepackage{pgfplots}
\usepackage{listings}
\usepackage{appendixnumberbeamer}

\pgfplotsset{compat=1.17}
\usetikzlibrary{positioning,arrows.meta,fit,backgrounds,calc,matrix,shapes.geometric,decorations.pathreplacing}

% =====================================================================
%  University of Helsinki palette
%  Primary brand colour is UH Hray/black with the "HY blue" accent.
%  Secondary accents chosen for projector contrast, not for print.
% =====================================================================
\definecolor{UHBlue}{HTML}{0C5A9E}   % primary accent
\definecolor{UHDark}{HTML}{15202B}   % body text / dark backgrounds
\definecolor{UHGrey}{HTML}{5A6672}   % secondary text
\definecolor{UHLight}{HTML}{EEF2F6}  % block backgrounds
\definecolor{UHOrange}{HTML}{D1601A} % alerts, gotchas
\definecolor{UHGreen}{HTML}{1F7A5C}  % examples, "this works"
\definecolor{UHRed}{HTML}{A62B2B}    % failure modes
\definecolor{UHSand}{HTML}{C8B88A}   % neutral third accent

\setbeamercolor{normal text}{fg=UHDark, bg=white}
\setbeamercolor{alerted text}{fg=UHOrange}
\setbeamercolor{example text}{fg=UHGreen}
\setbeamercolor{frametitle}{fg=white, bg=UHBlue}
\setbeamercolor{title separator}{fg=UHBlue}
\setbeamercolor{progress bar}{fg=UHBlue}
\setbeamercolor{progress bar in head/foot}{fg=UHBlue}
\setbeamercolor{progress bar in section page}{fg=UHBlue}
\setbeamercolor{structure}{fg=UHBlue}
\setbeamercolor{palette primary}{fg=white, bg=UHBlue}   % section pages, standout frames
\setbeamercolor{block title}{fg=white, bg=UHBlue}
\setbeamercolor{block body}{fg=UHDark, bg=UHLight}
\setbeamercolor{block title alerted}{fg=white, bg=UHOrange}
\setbeamercolor{block body alerted}{fg=UHDark, bg=UHLight}
\setbeamercolor{block title example}{fg=white, bg=UHGreen}
\setbeamercolor{block body example}{fg=UHDark, bg=UHLight}
\setbeamercolor{footnote}{fg=UHGrey}

\setbeamerfont{frametitle}{size=\large, series=\bfseries}
\setbeamerfont{footnote}{size=\tiny}
\setbeamerfont{caption}{size=\scriptsize}

% Frame title bar.
%
% Do NOT reimplement the frametitle template here. Metropolis's own template
% already spans \paperwidth (so `bg=UHBlue' above gives the full-width bar),
% already struts two-line titles without collision, and already draws the
% frametitle progress bar. Reimplementing it means touching internals whose
% names contain `@', which is NOT a letter in a .tex file -- `\@empty' parses
% as `\@' followed by the word "empty" and the build breaks.
%
% To change the bar's height, adjust Metropolis's own padding length instead,
% inside \makeatletter so that `@' is a letter:
%
%   \makeatletter
%   \setlength{\metropolis@frametitle@padding}{2.4ex}  % default 2.2ex
%   \makeatother
%
\makeatletter
\setlength{\metropolis@frametitle@padding}{1.2ex}
\makeatother

% No deck uses \framesubtitle; if one ever does, Metropolis ignores it and the
% subtitle should go in the frame body instead.

% =====================================================================
%  Course metadata — decks override \lecturename and \lecturenumber
% =====================================================================
\newcommand{\coursecode}{LDA-T318}
\newcommand{\coursename}{LLMs and Generative AI for NLP}
\newcommand{\courseteacher}{Aarne Talman}
\newcommand{\courseinstitution}{University of Helsinki}
\providecommand{\lecturename}{}
\providecommand{\lecturedate}{}

% Footline: course code on the left, lecture name centred, n/N on the right.
\setbeamertemplate{frame footer}{\coursecode\quad\textbar\quad\lecturename}

% =====================================================================
%  Semantic macros — use these instead of ad-hoc formatting, so that a
%  change of house style is a single edit here.
% =====================================================================

% --- Sourcing. Every quantitative claim on a slide should carry one. ---
\newcommand{\source}[1]{%
  \vfill{\scriptsize\color{UHGrey}\raggedright Source: #1\par}%
}
\newcommand{\sources}[1]{%
  \vfill{\scriptsize\color{UHGrey}\raggedright Sources: #1\par}%
}
% Inline citation chip, e.g. \cit{Raschka 2026a}
\newcommand{\cit}[1]{{\scriptsize\color{UHGrey}[#1]}}

% --- Counter-intuitive point. Three per lecture; these are exam material. ---
\newenvironment{gotcha}[1][Gotcha]{%
  \begin{alertblock}{#1}%
}{%
  \end{alertblock}%
}

% --- Speculative / unconfirmed claim. Research literacy matters here. ---
\newcommand{\rumour}{{\color{UHOrange}\textbf{[unconfirmed]}}\ }

% --- Key takeaway, for the closing slide of each section. ---
\newenvironment{takeaway}{%
  \begin{block}{Take away}%
}{%
  \end{block}%
}

% --- Live demo marker. ---
\newcommand{\demo}[1]{{\color{UHGreen}\textbf{Demo:}} #1}

% --- Emphasis consistent with the palette. ---
\newcommand{\key}[1]{{\color{UHBlue}\textbf{#1}}}
\newcommand{\warn}[1]{{\color{UHOrange}\textbf{#1}}}
\newcommand{\good}[1]{{\color{UHGreen}\textbf{#1}}}
\newcommand{\bad}[1]{{\color{UHRed}\textbf{#1}}}

% --- Units and quantities used constantly across the decks. ---
\newcommand{\unit}[1]{\,\mathrm{#1}}
\newcommand{\flops}{\mathrm{FLOPs}}
% NB: no \toks macro here --- \toks is a TeX primitive (a token register), so
% \newcommand{\toks} is a hard "already defined" error. No deck used it.
\newcommand{\KV}{\mathrm{KV}}
\newcommand{\softmax}{\operatorname{softmax}}

% =====================================================================
%  Code listings — used for config snippets, Slurm scripts, chat templates.
%  Frames containing these MUST be declared [fragile].
% =====================================================================
\lstdefinestyle{uh}{
  basicstyle=\ttfamily\scriptsize\color{UHDark},
  keywordstyle=\color{UHBlue}\bfseries,
  commentstyle=\color{UHGrey}\itshape,
  stringstyle=\color{UHGreen},
  numberstyle=\tiny\color{UHGrey},
  backgroundcolor=\color{UHLight},
  frame=none,
  xleftmargin=0.6em,
  framexleftmargin=0.6em,
  breaklines=true,
  breakatwhitespace=true,
  showstringspaces=false,
  columns=fullflexible,
  keepspaces=true,
  upquote=true,
  literate={-}{{-}}1 {_}{{\_}}1,
}
\lstset{style=uh}

\lstdefinelanguage{slurm}{
  morekeywords={SBATCH,module,srun,source,export,set,echo},
  morecomment=[l]{\#},
  sensitive=true,
}
\lstdefinelanguage{jinja}{
  morekeywords={for,endfor,if,endif,set,in},
  morecomment=[s]{\{\#}{\#\}},
  sensitive=true,
}

% =====================================================================
%  TikZ styles shared across decks (block diagrams, token grids).
% =====================================================================
\tikzset{
  mod/.style     = {draw=UHBlue, fill=UHLight, rounded corners=2pt,
                    minimum height=7mm, minimum width=22mm,
                    font=\scriptsize, align=center},
  modhi/.style   = {mod, draw=UHOrange, fill=UHOrange!12},
  modmute/.style = {mod, draw=UHGrey, fill=white, text=UHGrey},
  flow/.style    = {-{Latex[length=2mm]}, draw=UHGrey, thick},
  flowhi/.style  = {flow, draw=UHOrange},
  resid/.style   = {-{Latex[length=2mm]}, draw=UHBlue, thick, dashed},
  tokcell/.style = {draw=UHGrey!50, minimum size=4.2mm, font=\tiny, inner sep=0pt},
  tokon/.style   = {tokcell, fill=UHBlue!35},
  tokoff/.style  = {tokcell, fill=UHLight},
  tokmask/.style = {tokcell, fill=UHGrey!15, text=UHGrey!60},
  lbl/.style     = {font=\scriptsize\color{UHGrey}},
}

% =====================================================================
%  Presenter notes.
%
%  Notes are written on every slide. By default they are NOT shown, so
%  `latexmk -xelatex main.tex` produces a clean deck for projection.
%
%  To present with notes on a second screen (pdfpc, Skim, Impressive):
%      uncomment the line below, rebuild to main-notes.pdf
%  To print a notes handout:
%      \setbeameroption{show only notes}
% =====================================================================
% \setbeameroption{show notes on second screen=right}

\setbeamertemplate{note page}{%
  {\scriptsize\color{UHGrey}\coursecode\ \textbar\ \lecturename\par}%
  \vskip0.3em\hrule\vskip0.6em%
  {\small\insertnote}%
}

% ---------------------------------------------------------------- misc
\setbeamertemplate{caption}[numbered]
\setbeamertemplate{itemize items}[circle]
\setlength{\parskip}{0.1em}

% ---------------------------------------------------------------- density
% These decks are text-dense and a 16:9 frame is only 90mm tall, so the
% default vertical rhythm costs several lines per slide. Tightened globally
% here rather than slide by slide.
\setbeamersize{text margin left=0.7cm, text margin right=0.7cm}
\makeatletter
\def\@listi{\leftmargin\leftmargini
  \topsep 1.5\p@ \@plus1\p@ \@minus1\p@
  \parsep 0\p@
  \itemsep 1.5\p@ \@plus1\p@ \@minus1\p@}
\let\@listI\@listi
\makeatother

% The title page is a stack of short paragraphs, so a non-zero \parskip
% pushes it past the bottom of the frame. Neutralise it there only.
\addtobeamertemplate{title page}{\setlength{\parskip}{0pt}}{}

% Tighter lists — the decks are dense.
\setbeamertemplate{itemize/enumerate body begin}{\small}
\setbeamertemplate{itemize/enumerate subbody begin}{\footnotesize}

% Table helpers
\newcolumntype{L}[1]{>{\raggedright\arraybackslash}p{#1}}
\newcolumntype{C}[1]{>{\centering\arraybackslash}p{#1}}
\newcolumntype{R}[1]{>{\raggedleft\arraybackslash}p{#1}}

% Allow slightly overfull boxes without a warning storm
\hfuzz=2pt

% =====================================================================
%  Shared course-map frame.
%  Every lecture opens with this, with the current week highlighted:
%      \coursemapframe{1}   ... \coursemapframe{4}
%  Pass 0 to highlight nothing.
% =====================================================================
% \cmrow{current week}{this row's week}{dates}{Tuesday}{Thursday}{teacher}
% Highlights the row in UH blue bold when the two week numbers match.
% Deliberately avoids \rowcolor so that colortbl is not needed.
% The & and \\ must NOT sit inside the \ifnum branches: TeX scans alignment
% cells while skipping the false branch, so an & there breaks the conditional
% ("Incomplete \ifnum"). Keep the branches down to a macro name and let the
% separators live in the macro bodies, which expand after the test resolves.
\newcommand{\cmhi}[5]{%
  \textcolor{UHBlue}{\textbf{#1}} & \textcolor{UHBlue}{\textbf{#2}} &
  \textcolor{UHBlue}{\textbf{#3}} & \textcolor{UHBlue}{\textbf{#4}} &
  \textcolor{UHBlue}{\textbf{#5}}\\%
}
\newcommand{\cmplain}[5]{#1 & #2 & #3 & #4 & #5\\}
\newcommand{\cmrow}[6]{%
  \ifnum#1=#2\relax\expandafter\cmhi\else\expandafter\cmplain\fi
  {#2}{#3}{#4}{#5}{#6}%
}
\newcommand{\coursemapframe}[1]{%
\begin{frame}{Course overview}
\centering
\small
\setlength{\parskip}{0pt}
\renewcommand{\arraystretch}{1.08}
\begin{tabular}{@{}c L{1.7cm} L{3.1cm} L{3.7cm} L{1.4cm}@{}}
\toprule
\textbf{Wk} & \textbf{Dates} & \textbf{Tuesday} & \textbf{Thursday} & \textbf{Teacher}\\
\midrule
\cmrow{#1}{1}{27 / 29 Oct}{Language models}{Pretraining}{Aarne}
\cmrow{#1}{2}{3 / 5 Nov}{Post-training}{Lab 2: tune vs prompt}{Aarne}
\cmrow{#1}{3}{10 / 12 Nov}{Inference}{Lab 3: serve}{Aarne}
\cmrow{#1}{4}{17 / 19 Nov}{Evaluation}{Lab 4: evaluate}{Jussi}
\cmrow{#1}{5}{24 / 26 Nov}{RAG}{Lab 5: RAG}{Dmitry}
\cmrow{#1}{6}{1 / 3 Dec}{Agents}{Lab 6: agent}{Dmitry}
\cmrow{#1}{7}{8 / 10 Dec}{Course recap}{Special topics (TBC)}{All}
\bottomrule
\end{tabular}

\vspace{1.2em}
{\footnotesize Reusing a task and model across labs is recommended, not required.\\
There are no intermediate submissions; the final report and code repository
determine the grade.}
\note{Fifteen seconds, not more. The point of this slide is orientation, not
content: show them the shape of the course and where today sits inside it.
The final week is deliberately left flexible for a recap and special topics.}
\end{frame}%
}


\renewcommand{\lecturename}{Lecture 2 — Pretraining}
\renewcommand{\lecturedate}{29 October 2026}

\title{Pretraining}
\subtitle{Lecture 2: data, training and scale}
\author{Aarne Talman}
\institute{\coursecode\ \textbar\ \coursename\ \textbar\ \courseinstitution}
\date{29 October 2026}

\begin{document}

% =====================================================================
\begin{frame}[plain,shrink=8]
  \titlepage
  \note{On Tuesday we examined what a trained language model does. Today we
  look at how its weights are learned. The next-token loss is familiar; most
  of the new material concerns the data pipeline, compute budget and training
  system around it.

  This is the most quantitative lecture in the course. We derive one
  back-of-the-envelope equation, $C \approx 6ND$, and use it in Lab 1 and
  again in week 3.}
\end{frame}

\coursemapframe{1}

\begin{frame}{Lecture overview}
  \begin{columns}[T]
    \begin{column}{0.5\textwidth}
      \begin{enumerate}
        \item The training loop \hfill \textit{10 min}
        \item Data \hfill \textit{20 min}
        \item Scaling laws \hfill \textit{15 min}
        \item Training at scale \hfill \textit{15 min}
      \end{enumerate}
    \end{column}
    \begin{column}{0.5\textwidth}
      \begin{enumerate}
        \setcounter{enumi}{4}
        \item Mid-training \hfill \textit{10 min}
        \item What pretraining buys \hfill \textit{10 min}
        \item Lab 1 briefing \hfill \textit{8 min}
      \end{enumerate}
    \end{column}
  \end{columns}
  \vfill
  \begin{takeaway}
    By the end of today you should be able to explain how a language model is
    pretrained, estimate the compute and memory needed for a run, and identify
    the main choices in its data recipe.
  \end{takeaway}
  \note{Say the timings out loud again. Section 2 is the longest, deliberately
  so. Tell them that, because a twenty-minute block on data looks
  disproportionate to anyone who thinks this field is about architecture.

  The takeaway box is the promise. The three numbers are parameters, tokens
  and hardware-hours, and we will do exactly that calculation twice.}
\end{frame}

% =====================================================================
\section{The training loop}
% =====================================================================

\begin{frame}{Pretraining uses the same next-token objective}
  \[
    \mathcal{L}(\theta) \;=\; -\frac{1}{T}\sum_{t=1}^{T}
      \log p_\theta\!\left(x_t \mid x_{<t}\right)
  \]
  \emph{In words:} walk through the text one token at a time, ask how much
  probability the model gave to the token that actually came next, and average
  the penalty for being wrong. Same cross-entropy as Tuesday, now over a few
  trillion tokens.
  \vspace{0.3em}

  \key{Teacher forcing.} In training the model always conditions on the
  \emph{true} prefix, never on its own predictions. Every position is a
  supervised example, and the causal mask makes all of them computable in one
  forward pass.
  \vspace{0.3em}

  So a sequence of 4096 tokens provides a prediction target at almost every
  position, all computed in parallel.
  \vspace{0.4em}
  \small
  \textbf{Why this matters:} unlike an RNN, a transformer can process all
  training positions in parallel. Generation remains sequential; we return to
  that difference in week 3.
  \note{Open by writing the loss on the board and saying ``this is the same line
  as Tuesday''. Be almost boring about it. Students arrive expecting
  pretraining to involve some extra clever objective, and there is not one.

  Teacher forcing is the term they need for the reading. Make the point that it
  creates a train/test mismatch, since at inference the model conditions on its
  own output, and that the mismatch is one honest explanation for why long
  generations drift.

  The 4096-examples-per-forward-pass framing is the one that usually gets a
  visible reaction.}
\end{frame}

\begin{frame}[shrink=17]{Batches are measured in tokens}
  In LLM training, batch size is usually reported in \key{tokens}. The useful
  quantity is how many token positions contribute to one optimiser step.

  \vspace{0.5em}
  \centering
  \begin{tabular}{@{}L{3.0cm} L{3.2cm} L{4.6cm}@{}}
    \toprule
    \textbf{Run} & \textbf{Tokens / step} & \textbf{Note}\\
    \midrule
    GPT-3 175B (2020) & 3.2M & fixed through training\\[0.3em]
    Llama 3 405B (2024) & 4M \textrightarrow{} 8M \textrightarrow{} 16M & ramped, along with
      sequence length 4k \textrightarrow{} 8k\\[0.3em]
    Your Lab 1 run & 0.1M & one GPU, no accumulation needed\\
    \bottomrule
  \end{tabular}

  \vspace{0.8em}
  \raggedright
  \textbf{Gradient accumulation} buys you a large batch on small hardware:
  split the batch into several micro-batches, add up their gradients, and take
  one step. With the loss scaled correctly, this approximates one larger batch,
  subject to numerical and stochastic differences.
  \vspace{0.4em}

  \textbf{Batch ramping} starts small and grows. Early on the gradients are
  large, so a noisy estimate is good enough and a small batch gets more out of
  each token. Later the model needs the hardware efficiency of a big one.
  \source{Brown et al.\ 2020; Grattafiori et al., \textit{The Llama 3 Herd of Models}, 2024.}
  \note{The thing to land is the change of unit. If someone reports ``batch size
  64'' for an LLM they have told you almost nothing, because you do not know the
  sequence length.

  Gradient accumulation is worth thirty seconds. It is the most useful trick
  for anyone with one GPU, and they will use it in week 2.

  The Llama 3 ramp is a nice concrete: 4M to 16M tokens per step, changed twice
  during a 15T-token run, at 252M and 2.87T tokens.}
\end{frame}

\begin{frame}[shrink=12]{Why AdamW is commonly used}
  Plain gradient descent takes the same size of step for every weight. AdamW
  keeps \textbf{two running averages per weight} and uses them to size each
  weight's step individually.
  \vspace{0.4em}
  \begin{columns}[T]
    \begin{column}{0.53\textwidth}
      \key{Momentum} --- the average of this weight's recent gradients. Smooths
      out the noise from one batch to the next.
      \vspace{0.4em}

      \key{Second moment} --- the average of its recent gradients
      \emph{squared}: how big this weight's gradients typically are.
      \[
        \text{step} \;=\; \text{learning rate} \times
          \frac{\text{smoothed gradient}}{\text{typical gradient size}}
      \]
      \emph{In words:} divide each gradient by how large that weight's gradients
      usually are, so every weight moves by a comparable \emph{relative} amount
      whatever scale it lives on.
    \end{column}
    \begin{column}{0.47\textwidth}
      \begin{block}{What it costs}
        \small
        Those two averages are two extra full-size copies of the model, held in
        fp32.\\[0.4em]
        That is \textbf{8 bytes per parameter} of optimiser state, about half
        of the total training memory. It is the main reason a model you can
        serve on one card still will not train on it.
      \end{block}
      \vspace{0.3em}
      \small
      The \key{W} is decoupled weight decay: shrink every weight a little
      towards zero on each step, applied to the weights directly rather than
      pushed through the gradient.
    \end{column}
  \end{columns}
  \source{Loshchilov \& Hutter 2019; hyperparameters as used in GPT-3, Llama and
          most open 2024--2026 recipes.}
  \note{Do not derive Adam. Just insist on the two consequences: adaptivity is
  what makes a single learning rate work across very differently scaled
  parameter groups, and the two running averages cost eight bytes per
  parameter.

  The full update, if anyone wants it on the board. With $g_t$ the gradient at
  step $t$: $m_t = \beta_1 m_{t-1} + (1-\beta_1) g_t$ is the momentum buffer and
  $v_t = \beta_2 v_{t-1} + (1-\beta_2) g_t^2$ the second-moment buffer; after
  bias correction to $\hat m_t, \hat v_t$ the parameter update is
  $\theta_t = \theta_{t-1} - \eta(\hat m_t / (\sqrt{\hat v_t}+\epsilon) +
  \lambda\theta_{t-1})$.

  Why adaptivity matters: gradient magnitudes differ by orders of magnitude
  between embeddings, attention projections and norms, so one global learning
  rate cannot serve all of them. Dividing by the root second moment gives each
  parameter its own effective step size.

  Decoupled decay in full: the $\lambda\theta$ term acts on the parameters, not
  folded into the gradient. Folded in, Adam's normalisation would make the decay
  depend on gradient history, which nobody wants. That is the whole content of
  the 2019 paper.

  LLM hyperparameter convention: $\beta_1 = 0.9$, $\beta_2 = 0.95$ (not the
  PyTorch default $0.999$, because long-horizon second-moment estimates adapt too
  slowly on a trillion-token run), weight decay $\lambda = 0.1$.

  They need to remember the eight bytes. It reappears in twenty minutes in the
  memory accounting, and again in week 2 when we justify LoRA.}
\end{frame}

\begin{frame}{Keeping a training run stable}
  \begin{columns}[T]
    \begin{column}{0.5\textwidth}
      \textbf{Gradient clipping.} If the gradient as a whole is longer than a
      fixed threshold, almost always 1.0, scale it back down to that length.
      One bad batch then costs you a single bounded step rather than the run.
      \vspace{0.5em}

      \textbf{Warmup.} AdamW's running averages are poor estimates for the first
      few hundred steps, so start the learning rate near zero and ramp it up
      over roughly 0.1--2\% of training.
      \vspace{0.5em}

      \textbf{QK-Norm.} From Tuesday: normalise the queries and keys inside
      attention. Adopted mainly because it suppresses loss spikes.
    \end{column}
    \begin{column}{0.5\textwidth}
      \begin{block}{z-loss}
        \small
        A small extra penalty that keeps the model's raw output scores from
        drifting to ever-larger absolute values.
        \vspace{0.3em}
        That drift can build slowly and become expensive to correct late in a
        run. PaLM introduced this technique, and some later large-model recipes
        use it.
      \end{block}
      \vspace{0.3em}
      \small
      None of these make the model better. They are there so that the run
      reaches the end.
    \end{column}
  \end{columns}
  \source{Chowdhery et al., \textit{PaLM}, 2022; Raschka,
          \textit{A Dream of Spring for Open-Weight LLMs}, Feb 2026, on QK-Norm.}
  \note{Frame these as insurance. Each costs a little and buys a reduction in
  the probability of a catastrophic run failure, which on a three-month job is
  an excellent trade.

  Precise forms if asked. Clipping: rescale the gradient whenever its Euclidean
  norm $\|g\|_2$ exceeds $c$, with $c = 1.0$ standard. z-loss: add
  $10^{-4}\log^2 Z$ to the objective, where $Z$ is the softmax normaliser over
  the output vocabulary. Penalising $\log Z$ away from zero keeps the logits
  from inflating. QK-Norm is RMSNorm applied to $q$ and $k$ before the dot
  product.

  Say the last line on the slide out loud. Experienced people read a training
  recipe with that distinction in mind.}
\end{frame}

\begin{frame}{Learning-rate schedules: cosine versus WSD}
  \centering
  \begin{tikzpicture}
    \begin{axis}[
      width=10.4cm, height=4.4cm,
      xlabel={\scriptsize training progress (\% of planned tokens)},
      ylabel={\scriptsize LR / peak},
      xlabel style={font=\scriptsize}, ylabel style={font=\scriptsize},
      tick label style={font=\scriptsize},
      xmin=0, xmax=100, ymin=0, ymax=1.15,
      legend style={font=\scriptsize, draw=none, fill=none, at={(0.98,0.95)}, anchor=north east},
      axis lines=left, no markers, grid=none,
    ]
      \addplot[UHBlue, thick] coordinates {
        (0,0) (2,1) (10,0.98) (20,0.92) (30,0.83) (40,0.71) (50,0.58)
        (60,0.45) (70,0.32) (80,0.21) (90,0.13) (100,0.10)};
      \addlegendentry{warmup + cosine}
      \addplot[UHOrange, thick, dashed] coordinates {
        (0,0) (2,1) (20,1) (40,1) (60,1) (80,1) (85,0.72) (90,0.45)
        (95,0.20) (100,0.02)};
      \addlegendentry{warmup--stable--decay}
    \end{axis}
  \end{tikzpicture}
  \vspace{0.3em}

  \begin{columns}[T]
    \begin{column}{0.48\textwidth}
      \small
      \textbf{Cosine} needs the total token budget declared \emph{in advance}.
      Stop early and you stop at a high learning rate, with a model that is
      measurably worse than its loss curve suggests it should be.
    \end{column}
    \begin{column}{0.48\textwidth}
      \small
      \textbf{WSD} holds the peak rate flat, then decays sharply at the end.
      You can branch from \emph{any} point in the stable phase, anneal on a
      different mixture, and get a finished model.
    \end{column}
  \end{columns}
  \source{Hu et al., \textit{MiniCPM}, 2024; used in OLMo 2, SmolLM3 and
          several recent open-model recipes.}
  \note{WSD has become common partly for organisational reasons. A flat stable
  phase means one pretraining run can be
  forked into several finished models at different token budgets, and it means
  you do not have to commit to a total budget on day one, which you could not
  do anyway, because your cluster allocation will change.

  Point at the end of the orange curve and say: that final decay is where
  mid-training happens. We come back to it in section five.}
\end{frame}

\begin{frame}{Loss spikes}
  \centering
  \begin{tikzpicture}
    \begin{axis}[
      width=10.4cm, height=4.6cm,
      xlabel={\scriptsize tokens seen (billions)},
      ylabel={\scriptsize training loss},
      xlabel style={font=\scriptsize}, ylabel style={font=\scriptsize},
      tick label style={font=\scriptsize},
      xmin=0, xmax=100, ymin=2, ymax=11,
      axis lines=left, no markers,
    ]
      \addplot[UHBlue, thick] coordinates {
        (0,10.9) (1,7.2) (2,6.1) (4,5.3) (6,4.9) (8,4.6) (10,4.4)
        (14,4.1) (18,3.9) (22,3.75) (26,3.62) (30,3.52) (34,3.44)
        (36,3.40) (37,3.38) (38,6.9) (39,5.2) (40,4.3) (41,3.85)
        (42,3.60) (44,3.42) (46,3.34) (50,3.24) (55,3.14) (60,3.06)
        (65,2.99) (70,2.93) (75,2.88) (80,2.83) (85,2.79) (90,2.75)
        (95,2.72) (100,2.69)};
      \node[lbl, anchor=west] at (axis cs:42,8.0) {spike at 38B tokens};
      \draw[flow, draw=UHOrange] (axis cs:42,7.8) -- (axis cs:38.6,7.1);
    \end{axis}
  \end{tikzpicture}
  \vspace{0.2em}

  \begin{columns}[T]
    \begin{column}{0.49\textwidth}
      \small
      \textbf{Causes}: a corrupted or pathological data shard; learning rate
      too high for the current curvature; fp16 overflow; router collapse in an
      MoE; a silently faulty GPU.
    \end{column}
    \begin{column}{0.49\textwidth}
      \small
      \textbf{Mitigations}: QK-Norm and z-loss up front; skip the offending
      batches; roll back to a checkpoint about 100 steps earlier and skip
      200--500 batches of data. PaLM did exactly that, about twenty times.
    \end{column}
  \end{columns}
  \source{Chowdhery et al., \textit{PaLM}, 2022.}
  \note{The PaLM detail is the one worth telling in full, because it is so
  unglamorous. The team restarted from a checkpoint roughly a hundred steps
  before each spike and skipped a few hundred data batches. Re-running the same
  data from the same checkpoint reproduced the spike; skipping the data did not.
  So the data was the culprit rather than the optimiser, and they never found
  out which document.

  Ask the room what they would do if a spike recovers on its own. The answer is
  usually: nothing, but write down the step number, because the model quality
  after a recovered spike is not always the same.

  Clock check: should be around 16:25 here.}
\end{frame}

\begin{frame}{A healthy curve, and a comparison to avoid}
  \begin{columns}[T]
    \begin{column}{0.5\textwidth}
      \textbf{What good looks like in Lab 1}
      \begin{itemize}
        \item a near-vertical drop in the first few hundred steps, as the model
          learns the unigram distribution
        \item a long, smooth, slowly flattening descent on a log-x axis
        \item validation tracking training, with a gap that does not widen
        \item no early flat plateau; if the loss stops improving almost
          immediately, check the learning rate and data loader
      \end{itemize}
    \end{column}
    \begin{column}{0.5\textwidth}
      \begin{gotcha}[A lower loss is not always a better model]
        Cross-entropy is measured \textbf{per token}, so it is denominated in
        whatever your tokenizer calls a token.
        \vspace{0.3em}
        A larger vocabulary gives you fewer, harder tokens per sentence, so an
        equally good model records a \emph{higher} loss. Put more repetition in
        the mixture and the loss drops even though the model is worse.
        \vspace{0.3em}
        Loss only compares within one tokenizer and one data mixture. Across
        them, use \textbf{bits per byte}, or downstream accuracy.
      \end{gotcha}
    \end{column}
  \end{columns}
  \note{This gotcha is the direct continuation of Tuesday's perplexity gotcha
  and it is exam material. In Lab 1 they will train the same architecture with
  two or three vocabulary sizes and the temptation to rank them by validation
  loss is overwhelming. Tell them explicitly: that ranking is meaningless, and
  the lab asks for loss against compute at fixed tokenizer instead.

  Bits per byte is the fix worth naming: divide the total negative log
  likelihood by the number of UTF-8 bytes in the text rather than the number of
  tokens, and the unit is tokenizer-independent.

  The early-plateau warning is practical debugging advice. A run that barely
  improves usually has a configuration or data problem.}
\end{frame}

% =====================================================================
\section{Data}
% =====================================================================

\begin{frame}{The pipeline}
  \centering
  \begin{tikzpicture}[node distance=3mm]
    \node[mod, minimum width=11mm, minimum height=10mm, font=\tiny] (s1) {crawl};
    \node[mod, minimum width=11mm, minimum height=10mm, font=\tiny, right=of s1] (s2) {text\\extract};
    \node[mod, minimum width=11mm, minimum height=10mm, font=\tiny, right=of s2] (s3) {lang\\ID};
    \node[mod, minimum width=11mm, minimum height=10mm, font=\tiny, right=of s3] (s4) {dedup};
    \node[modhi, minimum width=11mm, minimum height=10mm, font=\tiny, right=of s4] (s5) {quality\\filter};
    \node[mod, minimum width=11mm, minimum height=10mm, font=\tiny, right=of s5] (s6) {PII};
    \node[mod, minimum width=11mm, minimum height=10mm, font=\tiny, right=of s6] (s7) {decon-\\taminate};
    \node[modhi, minimum width=11mm, minimum height=10mm, font=\tiny, right=of s7] (s8) {mixture};
    \draw[flow] (s1) -- (s2); \draw[flow] (s2) -- (s3); \draw[flow] (s3) -- (s4);
    \draw[flow] (s4) -- (s5); \draw[flow] (s5) -- (s6); \draw[flow] (s6) -- (s7);
    \draw[flow] (s7) -- (s8);
    \draw[decorate, decoration={brace, amplitude=4pt, mirror}, draw=UHGrey]
      (s1.south west) -- (s8.south east)
      node[midway, below=7pt, font=\scriptsize, color=UHGrey]
      {roughly $100$\,TB of HTML per CommonCrawl snapshot
       $\;\longrightarrow\;$ a few per cent survives};
  \end{tikzpicture}
  \vspace{1.2em}

  Each stage removes unwanted material, but it can also remove useful text.
  We will look at what each stage is for, where it can fail, and why the
  trade-offs differ across languages.
  \vfill
  \begin{takeaway}
    The data pipeline strongly affects what a model learns, yet model reports
    often document it less fully than the architecture.
  \end{takeaway}
  \note{Put the belt up and read the eight stages aloud. Tell them we spend
  twenty minutes here, a fifth of the lecture, and that the proportion is
  honest: this is where the quality comes from.

  The orange boxes are stages where model developers make especially
  consequential choices. The other stages still require careful engineering.

  Worth saying once: almost every frontier lab treats its data mixture as the
  crown jewel. OLMo and FineWeb are the honourable exceptions, which is exactly
  why our worked examples come from them.}
\end{frame}

\begin{frame}{Stage 1: the crawl}
  \begin{columns}[T]
    \begin{column}{0.54\textwidth}
      \key{CommonCrawl} is a non-profit that has crawled the web roughly
      monthly since 2008. Each snapshot is a few billion pages.
      \vspace{0.5em}

      It ships two formats, and the choice between them is the first quality
      decision in the pipeline:
      \begin{itemize}
        \item \key{WARC} --- the raw HTTP response. Full HTML, headers, the lot.
        \item \key{WET} --- CommonCrawl's own plain-text extraction, already done
          for you.
      \end{itemize}
      \vspace{0.5em}
      WET is convenient, but its extraction often leaves navigation bars,
      cookie banners, footers and link farms in the text.
    \end{column}
    \begin{column}{0.46\textwidth}
      \begin{block}{Failure mode of this stage}
        \small
        You inherit somebody else's definition of what a web page says.
        \vspace{0.4em}
        Every serious 2024--2026 dataset --- FineWeb, RefinedWeb,
        Dolma, HPLT --- starts from \textbf{WARC} and does its own extraction.
        \vspace{0.4em}
        The crawl is not neutral either. It respects \texttt{robots.txt}, and
        it under-samples paywalled and login-gated text while over-sampling
        whatever is cheap to serve.
      \end{block}
    \end{column}
  \end{columns}
  \source{CommonCrawl; Penedo et al., \textit{The FineWeb Datasets}, 2024;
          Penedo et al., \textit{RefinedWeb}, 2023.}
  \note{The WARC-versus-WET choice looks like plumbing and decides quality,
  which makes it the clearest example of its kind in the lecture. Teams that
  used WET in 2021 and 2022 were effectively training on boilerplate.

  Mention the sampling bias at the end. It is what makes the web a strange
  corpus: it over-represents what is cheap and public, which is not the same
  thing as what is good.}
\end{frame}

\begin{frame}{Stage 2: text extraction}
  Turning HTML into the article, and only the article.
  \vspace{0.6em}
  \begin{columns}[T]
    \begin{column}{0.52\textwidth}
      \textbf{What must go}: navigation, sidebars, cookie notices, ``related
      articles'', comment spam, licence footers, JavaScript.
      \vspace{0.5em}

      \textbf{What must stay}: the body text, in reading order, with paragraph
      and list structure intact.
      \vspace{0.5em}

      \key{trafilatura} is the standard tool. FineWeb measured the swap
      directly: WET $\to$ trafilatura on WARC improved downstream benchmark
      scores on models trained on the output, at the cost of a slower and much
      more expensive extraction step.
    \end{column}
    \begin{column}{0.48\textwidth}
      \begin{block}{Failure mode of this stage}
        \small
        Extraction errors come in \textbf{systematic} patterns rather than at
        random.
        \vspace{0.4em}
        A template that strips the first paragraph of every page on a large CMS
        removes exactly the lead sentences, the densest text on the page,
        across millions of documents at once.
        \vspace{0.4em}
        You will not notice this by sampling, because each individual document
        still looks fine.
      \end{block}
    \end{column}
  \end{columns}
  \source{Penedo et al., \textit{The FineWeb Datasets}, 2024; Barbaresi,
          \textit{trafilatura}, 2021.}
  \note{The failure mode is the point of the slide, and it foreshadows the
  gotcha four slides from now. Extraction bugs correlate with site templates,
  so they hit millions of documents in the same way, and the damaged documents
  still read as plausible text.

  The one-sentence version: at scale you can survive random corruption, but
  systematic corruption follows the model all the way through.}
\end{frame}

\begin{frame}{Stage 3: language identification}
  A classifier, usually fastText trained on Wikipedia and web text, gives each
  document a language and a confidence. You keep what clears a threshold.

  \vspace{0.5em}
  \centering
  \begin{tabular}{@{}L{3.2cm} L{7.8cm}@{}}
    \toprule
    \textbf{Where it breaks} & \textbf{Consequence}\\
    \midrule
    short documents & unreliable under about 20 words, so short texts vanish\\[0.25em]
    code-switching & a Finnish post quoting English clears neither threshold\\[0.25em]
    related languages & Danish/Bokm\aa{}l and Czech/Slovak get confused\\[0.25em]
    the threshold itself & FineWeb keeps English above a confidence of 0.65\\
    \bottomrule
  \end{tabular}

  \vspace{0.6em}
  \raggedright
  \small
  Each threshold trades quality against quantity, and a smaller language cannot
  give up quantity.
  \source{Joulin et al., fastText, 2016; Penedo et al., \textit{The FineWeb Datasets}, 2024.}
  \note{The threshold is the teaching point. Raise it and you lose data; lower
  it and you take on noise. For English you can afford a confidence of 0.65
  because there is more English than anyone can use. For Finnish, every point
  of threshold you raise is data you cannot replace, so Finnish pipelines run
  looser thresholds and then pay for it in noise.

  The code-switching row is worth a beat for this audience: multilingual
  speakers write mixed text constantly, and the standard pipeline treats that
  as garbage.

  Two rows that did not fit on the slide. Indonesian/Malay is another pair that
  language ID routinely confuses. And code and maths get classified as whatever
  natural language the comments happen to be in, or else as noise, which is one
  reason code is usually sourced separately from GitHub rather than pulled out
  of a web crawl.}
\end{frame}

\begin{frame}[shrink=7]{Stage 4: deduplication}
  \begin{columns}[T]
    \begin{column}{0.52\textwidth}
      Three levels, usually all three:
      \begin{itemize}
        \item \key{exact} --- hash the document, drop repeats. Cheap, catches
          mirrors and reposts.
        \item \key{near-duplicate} --- hash overlapping word sequences and
          compare the hashes (MinHash and LSH), dropping pairs that share about
          80\% of them. Catches templated and lightly edited pages.
        \item \key{line-level} --- drop lines repeated across many documents.
          This is what removes the surviving boilerplate.
      \end{itemize}
      \vspace{0.4em}
      Why it helps: duplicated text is effectively upweighted, which inflates
      memorisation and leaks test sets into training.
    \end{column}
    \begin{column}{0.48\textwidth}
      \begin{block}{The FineWeb surprise}
        \small
        Deduplicating \emph{globally} across all 96 CommonCrawl snapshots gave
        \textbf{worse} models than deduplicating within each snapshot.
        \vspace{0.4em}
        Why: text that survives on the web across many years and many sites is
        disproportionately \emph{good}. Aggressive global dedup removes it and
        leaves a residue of unique, never-copied, low-quality pages.
        \vspace{0.4em}
        So more filtering does not always give you a better corpus.
      \end{block}
    \end{column}
  \end{columns}
  \source{Lee et al., \textit{Deduplicating Training Data Makes LMs Better}, 2022;
          Penedo et al., \textit{The FineWeb Datasets}, 2024.}
  \note{Dedup is the one stage where the literature is unambiguous that it
  helps. Then FineWeb found the limit of that claim, which is why this slide is
  worth more than its share of time.

  Spell out the mechanism of the surprise: duplication on the web is a weak
  quality signal, because people copy things that are worth copying. Global
  dedup discards that signal entirely.

  Connect it forward: this is the first of three places today where the
  intuitive move turns out to be wrong at scale.}
\end{frame}

\begin{frame}[shrink=7]{Stage 5a: quality filtering by hand}
  Hand-written heuristics, inherited largely from C4 and Gopher. A document
  survives if, roughly:
  \vspace{0.5em}
  \begin{columns}[T]
    \begin{column}{0.5\textwidth}
      \small
      \begin{itemize}
        \item 50--100\,000 words, mean word length 3--10 characters
        \item fewer than 30\% of lines ending in an ellipsis
        \item at least 80\% of words containing an alphabetic character
        \item a sentence-ending punctuation mark on most lines
        \item symbol-to-word ratio below $0.1$
        \item no match against a list of blocked terms
      \end{itemize}
    \end{column}
    \begin{column}{0.5\textwidth}
      \begin{block}{What these rules really are}
        \small
        A set of cheap proxies for ``written by a person, for a reader''.
        \vspace{0.4em}
        They are also \textbf{English-shaped}. A mean word length of 3--10
        characters is wrong for Finnish or German compounding, and
        sentence-final punctuation is wrong for Thai. Blocked-term lists remove
        sex education along with pornography, and they take out a lot of text
        written by and about minority groups.
      \end{block}
    \end{column}
  \end{columns}
  \vspace{0.5em}
  \small
  The rules are cheap and you can audit them. They also put a ceiling on
  quality, because a rule you can write down is a rule a spammer can satisfy.
  \source{Raffel et al., \textit{C4}, 2020; Rae et al., \textit{Gopher}, 2021;
          Dodge et al., \textit{Documenting the English C4}, 2021.}
  \note{Read two or three of the rules aloud in a deadpan voice. Putting
  ``trillion-parameter frontier model'' next to ``fewer than thirty per cent of
  lines ending in an ellipsis'' is genuinely funny, and it makes the point
  stick.

  The fairness objection in the right-hand block is documented, not
  speculative: Dodge and colleagues audited C4's blocklist and found it removed
  a large amount of non-offensive content, with a measurable skew.

  The last line sets up the next slide.}
\end{frame}

\begin{frame}{Stage 5b: quality filtering by classifier}
  Replace the rules with a model. \key{FineWeb-Edu} is the canonical worked
  example.
  \vspace{0.5em}
  \begin{columns}[T]
    \begin{column}{0.56\textwidth}
      \begin{enumerate}
        \item Ask a large instruct model to rate 460k FineWeb documents 0--5
          for \emph{educational value}
        \item Train a small, cheap classifier on those labels, on top of an
          off-the-shelf embedding model
        \item Run it over all 15T tokens of FineWeb
        \item Keep the documents scoring 3 or above
      \end{enumerate}
      \vspace{0.5em}
      Result: \textbf{1.3T tokens} out of 15T, so under a tenth survives. At
      the same token budget, models trained on the survivors are substantially
      better on knowledge-and-reasoning benchmarks such as MMLU and ARC.
    \end{column}
    \begin{column}{0.44\textwidth}
      \begin{block}{The cost of the trick}
        \small
        You have encoded one model's taste into your corpus.
        \vspace{0.4em}
        ``Educational value'' favours textbooks, tutorials and encyclopaedic
        prose, and it disfavours dialogue, fiction, argument and informal
        registers.
        \vspace{0.4em}
        If you want a model that can write, that is the wrong thing to filter
        on, and you will discover it after the run.
      \end{block}
    \end{column}
  \end{columns}
  \source{Lozhkov et al., \textit{FineWeb-Edu}, 2024; Penedo et al.,
          \textit{The FineWeb Datasets}, 2024.}
  \note{The numbers to say: 460 thousand annotations from a big model, one
  small classifier, fifteen trillion tokens in, one point three trillion out.
  A ninety-one per cent discard rate, and the result is better.

  Then say the honest part. FineWeb-Edu is a benchmark-shaped dataset, and the
  FineWeb team say so themselves. Filter for a benchmark and you end up with a
  corpus shaped like that benchmark. If your downstream use is conversational
  or creative, the educational filter is actively the wrong objective.

  This is also the best available answer to ``can I just use an LLM to clean my
  data?'': yes, as an annotator for a cheap classifier, and no, not as the
  filter itself. Nobody can afford to run a 70B model over fifteen trillion
  tokens.}
\end{frame}

\begin{frame}[shrink=9]{How to compare data pipelines}
  \begin{gotcha}[Inspection is necessary, but not sufficient]
    Samples and corpus statistics can reveal obvious problems. To determine
    which pipeline produces better training data, train comparable small models
    and evaluate them on the same tasks.
  \end{gotcha}
  \vspace{0.5em}
  \begin{columns}[T]
    \begin{column}{0.54\textwidth}
      \textbf{The FineWeb ablation protocol}
      \begin{itemize}
        \item fix the architecture, tokenizer, hyperparameters and token budget
        \item train a 1.8B-parameter model on 28B tokens from each candidate
          dataset
        \item evaluate on a fixed suite of benchmarks chosen in advance for
          low variance and a signal above chance at that scale
        \item change \textbf{one} pipeline decision at a time
      \end{itemize}
    \end{column}
    \begin{column}{0.46\textwidth}
      \begin{block}{Why this is the methodological lesson}
        \small
        It turns data engineering into something you can run experiments on.
        \vspace{0.4em}
        It is also expensive: every filter you want to justify costs a training
        run, which is why the published ablation studies come from the handful
        of teams willing to spend compute on questions rather than on models.
        \vspace{0.4em}
        The same protocol is what Lab 1 Part C asks you to run, in miniature.
      \end{block}
    \end{column}
  \end{columns}
  \source{Penedo et al., \textit{The FineWeb Datasets}, 2024.}
  \note{This slide transfers further than anything else in the lecture, and it
  is exam material. Say the gotcha, pause, and then ask the room how they would
  decide between two candidate corpora. Someone will say ``look at them'',
  which is what almost everyone says and it is wrong.

  The benchmark-selection detail is worth ten seconds: you cannot use MMLU in a
  1.8B ablation because it is at chance, so you choose a suite that has signal
  at your ablation scale. Choosing the wrong suite makes the whole protocol
  blind.

  Clock check: should be around 16:40 here.}
\end{frame}

% Useful background, omitted from the 90-minute projected deck.
\iffalse
\begin{frame}{Stage 6: PII removal}
  \begin{columns}[T]
    \begin{column}{0.54\textwidth}
      Regular expressions and named-entity models over email addresses, phone
      numbers, IP addresses, and sometimes postal addresses and identity
      numbers. Matches are replaced with placeholders rather than deleted, so
      the surrounding text stays grammatical.
      \vspace{0.6em}

      There are two separate motivations here, and they often get conflated:
      \begin{itemize}
        \item \key{legal} --- personal data requires a lawful basis, data
          minimisation and appropriate risk controls under GDPR
        \item \key{behavioural} --- a model that has memorised contact details
          will emit them, and that is a product failure before it is a legal one
      \end{itemize}
    \end{column}
    \begin{column}{0.46\textwidth}
      \begin{block}{Failure mode of this stage}
        \small
        These tools miss a great deal, and they also flag a great deal they
        should not.
        \vspace{0.4em}
        Regexes miss anything obfuscated and destroy legitimate text: version
        numbers that look like IP addresses, Finnish place names that a NER
        model reads as person names.
        \vspace{0.4em}
        And stripping PII is not the same as privacy. A model can reconstruct
        an identity from context long after every email address has gone.
      \end{block}
    \end{column}
  \end{columns}
  \source{Soldaini et al., \textit{Dolma}, 2024; FineWeb-PII, 2024.}
  \note{Keep this short, but do not skip it. It is the stage with the clearest
  legal exposure for anyone in this room who ends up building a corpus in
  Europe.

  The line to land is the last one. PII scrubbing mitigates the problem and
  does not solve it, so anyone who claims their model cannot leak personal data
  because they ran a regex is overclaiming.}
\end{frame}
\fi

\begin{frame}[shrink=9]{Stage 7: decontamination}
  Remove from the training corpus any text that overlaps the \emph{test sets}
  you intend to report.
  \vspace{0.6em}
  \begin{columns}[T]
    \begin{column}{0.52\textwidth}
      \textbf{How}: $n$-gram overlap against every benchmark's test split.
      GPT-3 used 13-grams; later work uses 8- to 13-grams or normalised
      substring matching. Flagged documents are dropped, or the matching
      passage is excised.
      \vspace{0.5em}

      \textbf{Why it is hard}
      \begin{itemize}
        \item you must know your benchmarks before you train
        \item paraphrases and translations are invisible to $n$-gram matching
        \item a later benchmark can still overlap with earlier source material,
          paraphrases, translations or later training stages
      \end{itemize}
    \end{column}
    \begin{column}{0.48\textwidth}
      \begin{block}{Why you should care in week 4}
        \small
        Contamination can inflate a benchmark score without an obvious warning
        in the output.
        \vspace{0.4em}
        It is one possible explanation for an unexpectedly large result, and
        one reason Jussi's evaluation lecture emphasises held-out data you
        control.
        \vspace{0.4em}
        In your own project: your held-out set must never have been used during
        development, for exactly the same reason at a smaller scale.
      \end{block}
    \end{column}
  \end{columns}
  \source{Brown et al.\ 2020; Soldaini et al., \textit{Dolma}, 2024;
          Raschka, \textit{Understanding the 4 Main Approaches to LLM Evaluation}, 2025.}
  \note{Decontamination is the bridge to week 4 and to their own experimental
  hygiene, so make that link explicit rather than treating it as a data-cleaning
  footnote.

  A benchmark newer than the main crawl reduces one contamination route, but
  it is not a guarantee: source material, later training stages and semantic
  overlap still matter.}
\end{frame}

\begin{frame}[shrink=12]{Stage 8: the mixture}
  The last decision is how much of each source goes in, and how many times each
  one is seen. It has more effect on the model than anything before it.

  \vspace{0.4em}
  \centering
  \begin{tabular}{@{}L{2.8cm} L{2.4cm} L{5.6cm}@{}}
    \toprule
    \textbf{Source} & \textbf{Typical share} & \textbf{What it contributes}\\
    \midrule
    filtered web & 60--85\% & breadth, world knowledge, register variety\\[0.25em]
    code & 5--20\% & reasoning in general, not only coding\\[0.25em]
    maths & 2--10\% & arithmetic and multi-step derivation\\[0.25em]
    books, papers & 2--10\% & long-range coherence, formal register\\[0.25em]
    curated reference & 1--5\% & tiny, and upsampled aggressively\\
    \bottomrule
  \end{tabular}

  \vspace{0.5em}
  \raggedright
  \small
  \key{Upsampling} means sampling a scarce source more often. One study found
  useful gains over several passes and diminishing returns with heavy
  repetition; the useful limit depends on the corpus and model.
  \source{Muennighoff et al., \textit{Scaling Data-Constrained Language Models}, 2023;
          Aryabumi et al., \textit{To Code or Not To Code?}, 2024.}
  \note{Two numbers to land. First, code in the pretraining mixture improves
  general reasoning as well as code generation. Cohere ran that ablation
  properly and it is one of the more surprising robust findings in the field.
  Nobody fully agrees on why; the leading guess is that code supplies long,
  structured, verifiable dependency chains.

  The curated-reference row is mostly Wikipedia: a few gigabytes of text that
  gets seen many times over a trillion-token run.

  Second, the four-epoch rule. Anyone working in a data-constrained language
  will use that number constantly, and we come back to it on the Finnish slide
  in two minutes.

  Shares are indicative and vary a lot between labs. Say that out loud so
  nobody quotes the table as gospel.}
\end{frame}

\begin{frame}{A smaller language has a smaller data budget}
  \begin{columns}[T]
    \begin{column}{0.5\textwidth}
      \begin{itemize}
        \item English is roughly \textbf{45\%} of CommonCrawl documents;
          Finnish roughly \textbf{0.2\%}
        \item after filtering, usable Finnish web text is on the order of
          \textbf{tens of billions} of tokens
      \end{itemize}
      \vspace{0.4em}

      A compute-optimal 7B model uses about 140B tokens, and many recent
      recipes use substantially more. Finnish data alone is therefore not
      enough for this kind of general-purpose model.
    \end{column}
    \begin{column}{0.5\textwidth}
      \begin{block}{The consequence}
        \small
        A practical approach is a \textbf{multilingual} model in which Finnish
        is upsampled.
        \vspace{0.4em}
        This supplies enough varied training text while giving Finnish a larger
        share than its raw web frequency.
        \vspace{0.4em}
        It can also transfer useful patterns learned from English and code to
        Finnish contexts.
      \end{block}
    \end{column}
  \end{columns}
  \source{CommonCrawl language statistics; TurkuNLP Finnish web corpora;
          Luukkonen et al., \textit{Poro 34B}, 2024.}
  \note{This is where the lecture becomes local. Most of the room either speaks
  Finnish or works on a language with a similar data budget, so make the
  arithmetic concrete and let it land.

  The number that does the work is the ratio: two hundred to one against
  English before filtering, and the gap widens after filtering because the
  quality heuristics are English-shaped.

  The upside at the bottom is easy to miss: a model can carry reasoning ability
  from English into Finnish even though it cannot carry Finnish facts it never
  saw. That is the entire basis of the Poro design on the next slide.}
\end{frame}

\begin{frame}{Ways to increase training data for a smaller language}
  \small
  \centering
  \begin{tabular}{@{}L{2.7cm} L{4.2cm} L{4.2cm}@{}}
    \toprule
    \textbf{Approach} & \textbf{What it provides} & \textbf{Main limitation}\\
    \midrule
    multilingual mixture & more total data and cross-lingual transfer &
      less compute spent directly on the target language\\[0.35em]
    repeat scarce data & greater weight for existing target-language text &
      diminishing returns and more memorisation\\[0.35em]
    machine translation & large volumes from high-resource languages &
      translation patterns and MT errors\\[0.35em]
    synthetic text & controllable topics and formats &
      generator bias, errors and licence constraints\\[0.35em]
    parallel text & explicit links between languages &
      limited volume and domain coverage\\
    \bottomrule
  \end{tabular}

  \vspace{0.6em}
  \raggedright
  These approaches are usually combined. Poro, for example, mixed Finnish with
  English and code, repeated Finnish data, and included aligned
  English--Finnish text.
  \source{Muennighoff et al.\ 2023; Luukkonen et al., \textit{Poro 34B}, 2024.}
  \note{Walk the table quickly and stress that these are complementary rather
  than mutually exclusive choices. Parallel data is particularly relevant to
  the Poro case study.

  Translationese deserves a sentence: translated text has systematically
  different distributional properties from natively written text, and a model
  trained heavily on it writes like a translation. Several national-language
  model projects have run into this.}
\end{frame}

\begin{frame}[shrink=8]{Case study: Poro and Viking on LUMI}
  \begin{columns}[T]
    \begin{column}{0.52\textwidth}
      \textbf{Poro 34B} (2024) --- Silo AI with TurkuNLP and HPLT, trained on
      LUMI on roughly \textbf{1T tokens}:
      \begin{itemize}\small
        \item English (SlimPajama), the majority of the mixture
        \item \emph{all} the usable Finnish there was --- on the order of 30B
          tokens, repeated several epochs
        \item a large code share (StarCoder)
        \item explicit \key{English--Finnish translation pairs}
      \end{itemize}
      \vspace{0.4em}
      \textbf{Viking 7B / 13B / 33B} --- the follow-up: about 2T tokens
      across Finnish, English, Swedish, Danish, Norwegian, Icelandic and code.
    \end{column}
    \begin{column}{0.48\textwidth}
      \begin{block}{What it bought}
        \small
        Finnish performance far beyond anything trainable on Finnish alone,
        and English did not suffer for it. That is the cross-lingual transfer
        argument, demonstrated at scale on a national supercomputer.
      \end{block}
      \begin{block}{What I would do differently}
        \small
        The architecture was frozen roughly nine months before release, and the
        field moved. A token budget that looked ambitious at kick-off was
        unremarkable at launch. And a large share of the engineering effort
        went into the software stack rather than the model.
      \end{block}
    \end{column}
  \end{columns}
  \source{Luukkonen et al., \textit{Poro 34B and the Blessing of Multilinguality}, 2024;
          Viking model cards, Silo AI / TurkuNLP / HPLT, 2024.
          Mixture shares approximate.}
  \note{This is first-hand, so tell it as a story rather than reading the slide.
  The decision that is worth the most to them: we did not try to build a Finnish
  model, we built an English-and-code model that also had Finnish in it, because
  that was the only way the data budget closed.

  Be honest about the retrospective. The nine-month architecture freeze is the
  real lesson of large-scale training: the run usually finishes, and by the
  time it does the field has moved.

  If asked about compute: a few hundred LUMI GPUs for several weeks. Exact
  figures are in the paper and I would rather they look them up than that I
  misquote them.}
\end{frame}

\begin{frame}[shrink=11]{Licensing, and the open-data question}
  \begin{columns}[T]
    \begin{column}{0.52\textwidth}
      \textbf{What is actually released}
      \begin{itemize}
        \item \key{weights only} --- almost everything called ``open'':
          Llama, Qwen, Gemma, DeepSeek, Mistral
        \item \key{weights + data + code + recipe} --- OLMo, and very little else
        \item \key{data only} --- FineWeb, Dolma, HPLT, The Stack
      \end{itemize}
      \vspace{0.5em}
      The legal status of training on copyrighted web text is unsettled and
      differs between the EU and the US. The EU's text-and-data-mining exception
      has an opt-out that publishers increasingly exercise.
    \end{column}
    \begin{column}{0.48\textwidth}
      \begin{block}{Why this matters for research}
        \small
        If the data is secret, the experiment is not reproducible, and claims
        about architecture cannot be separated from claims about data.
        \vspace{0.4em}
        That is the concrete reason I said on Tuesday that ``which architecture
        is best'' is unanswerable from public evidence.
      \end{block}
      \vspace{0.3em}
      \small
      For your project: name the licence of every model and dataset you use, and
      check that it permits academic use. This is in the rubric.
    \end{column}
  \end{columns}
  \vfill
  \begin{takeaway}
    At the same scale, what separates two models is mostly what they were
    trained on, and that is the part of a recipe least likely to be documented.
  \end{takeaway}
  \note{Close the section on the research-literacy point, which is the one that
  connects back to Tuesday's honest closer about architecture.

  The OLMo exception is worth naming every time it comes up. It is the only
  family where you can actually run the controlled experiment, and that is why
  it is over-represented in the citations on these slides relative to its
  benchmark scores.

  Clock check: should be around 16:47 here. If we are late, cut the licensing
  discussion to thirty seconds. The scaling-law section cannot be compressed.}
\end{frame}

% =====================================================================
\section{Scaling laws}
% =====================================================================

\begin{frame}{Loss is a power law in compute}
  \centering
  \begin{tikzpicture}
    \begin{axis}[
      width=10.2cm, height=4.4cm,
      xmode=log, ymode=log,
      xlabel={\scriptsize training compute $C$ (FLOPs)},
      ylabel={\scriptsize loss},
      xlabel style={font=\scriptsize}, ylabel style={font=\scriptsize},
      tick label style={font=\scriptsize},
      xmin=3e16, xmax=3e25, ymin=1.4, ymax=5.5,
      axis lines=left, no markers,
      ytick={1.5,2,3,4,5}, yticklabels={1.5,2,3,4,5},
    ]
      \addplot[UHBlue, thick] coordinates {
        (1e17,4.60) (1e18,4.03) (1e19,3.52) (1e20,3.08) (1e21,2.70)
        (1e22,2.36) (1e23,2.06) (1e24,1.81) (1e25,1.58)};
      \addplot[only marks, mark=*, mark size=1.6pt, UHOrange] coordinates {
        (1.2e17,4.55) (7.2e23,1.90) (3.3e24,1.73)};
      \node[lbl, anchor=west, font=\tiny\color{UHGrey}] at (axis cs:1.6e17,4.0) {your Lab 1 run};
      \node[lbl, anchor=east, font=\tiny\color{UHGrey}] at (axis cs:5.0e23,1.62) {Llama 3 8B};
      \node[lbl, anchor=west, font=\tiny\color{UHGrey}] at (axis cs:4.0e24,1.95) {DeepSeek V3};
    \end{axis}
  \end{tikzpicture}

  \begin{columns}[T]
    \begin{column}{0.52\textwidth}
      \small
      Kaplan et al.\ (2020): across seven orders of magnitude, test loss falls
      approximately as a \key{power law} in parameters, data and compute within
      the range they measured. On log--log axes this appears close to a
      straight line.
    \end{column}
    \begin{column}{0.48\textwidth}
      \small
      This is what made scaling an \emph{engineering} activity. You can fit the
      line at small scale and predict the loss of a run you have not yet paid
      for. Frontier labs now do exactly that before committing.
    \end{column}
  \end{columns}
  \source{Kaplan et al., \textit{Scaling Laws for Neural Language Models}, 2020.
          Curve is illustrative; the markers use $C = 6ND$ from the next slides.}
  \note{Larger models were already known to help. The striking observation was
  how closely the measured losses followed smooth trends across a wide compute
  range.

  The practical consequence is the one to emphasise. Scaling laws let teams
  use smaller experiments to inform the budget and expected loss of a larger
  run.

  Flag that this curve is illustrative, since a real one is fitted from dozens
  of runs. The three markers, though, are placed with the arithmetic we are
  about to do.}
\end{frame}

\begin{frame}[shrink=5]{Chinchilla: how to split a fixed budget}
  You have a fixed compute budget. Do you spend it on a \textbf{bigger model}
  ($N$ parameters) or on \textbf{more text} ($D$ tokens)?
  \vspace{0.3em}

  Hoffmann et al.\ fitted the loss as a function of both, then asked which split
  minimises it with the budget $C = 6ND$ held fixed. The answer:
  \[
    \frac{D}{N} \;\approx\; 20\ \text{tokens per parameter}
  \]
  \emph{In words:} in their experiments, model size and data size should grow
  \key{together}. The value 20 is a useful historical reference, not a
  universal constant.
  \vspace{0.3em}
  \begin{columns}[T]
    \begin{column}{0.52\textwidth}
      \small
      \textbf{The demonstration.} Gopher: 280B parameters on 300B tokens.
      Chinchilla: \textbf{70B} parameters on \textbf{1.4T} tokens. Same
      training compute, but a smaller model trained on much more data.
      Chinchilla reported lower loss and stronger results on many benchmarks.
    \end{column}
    \begin{column}{0.48\textwidth}
      \small
      \textbf{How to use the result.} Treat the ratio as a starting point.
      Architecture, data quality, optimiser and learning-rate schedule all
      affect the fitted optimum, so current projects run their own small-scale
      studies.
    \end{column}
  \end{columns}
  \source{Hoffmann et al., \textit{Training Compute-Optimal Large Language Models}, 2022.}
  \note{The room may have heard ``twenty tokens per parameter'' without knowing
  where it comes from. It is the optimum fitted in this particular study, not
  a general law.

  The fitted form, for anyone who wants it:
  $L(N,D) \approx E + A/N^{\alpha} + B/D^{\beta}$ with $\alpha \approx 0.34$ and
  $\beta \approx 0.28$. $E$ is the irreducible entropy of the text, the second
  term is the cost of a finite model, the third the cost of finite data.
  Minimising this subject to $C = 6ND$ gives $N \propto C^{1/2}$ and
  $D \propto C^{1/2}$, so both grow as the square root of compute, and the
  ratio $D/N \approx 20$. The reason both exponents come out at one half is that
  $\alpha$ and $\beta$ are nearly equal; had one been much larger, the optimum
  would have been lopsided.

  The Gopher comparison is the headline: the previous generation had been
  spending its compute in the wrong place for years, and a model four times
  smaller beat it on the same budget. Everything got cheaper to serve
  overnight.

  The research-literacy point is that scaling-law coefficients are empirical.
  They depend on the models, data and training procedure used to fit them.}
\end{frame}

\begin{frame}{Where $6ND$ comes from}
  A \key{FLOP} is one floating-point operation: a single multiply or a single
  add. Using one weight on one token means one multiply and one add, so
  \textbf{2 FLOPs per weight per token}.
  \vspace{0.5em}
  \begin{itemize}
    \item \textbf{Forward pass.} Every token passes through every parameter
      once. \hfill $2ND$
      \vspace{0.4em}
    \item \textbf{Backward pass.} Twice the work of the forward pass, because
      each layer computes two gradients: one to pass further back down the
      network, one to update its own weights. \hfill $4ND$
  \end{itemize}
  \vspace{0.5em}
  \[ \boxed{\,C \;\approx\; 6 N D\,} \qquad
     N = \text{parameters}, \quad D = \text{training tokens} \]
  \emph{In words:} a training run costs about six operations per parameter per
  token for the parameterised layers of a conventional dense transformer. The
  next slide lists important costs this estimate omits.
  \note{Do this on the board, step by step, and do not skip the factor of two
  from multiply-plus-add. If you assert it, half the class loses the thread
  there.

  The shape argument, if you want it in full: multiplying a $B \times m$ block of
  activations by an $m \times n$ weight matrix costs $2Bmn$ FLOPs, which is two
  per input element per weight. Summed over all weights and all tokens that is
  $2ND$ forward.

  The backward pass is the step people forget: it is two matrix multiplies of
  the same shape per layer, not one, because you need the gradient with respect
  to the layer's \emph{input} in order to propagate, and the gradient with
  respect to its \emph{weights} in order to update. So $2 \times 2ND = 4ND$, and
  the total is $6ND$. Draw the chain rule for a single linear layer if anyone
  looks unconvinced.

  Then say the point of the whole slide: they own this number now, and can
  rebuild it in thirty seconds, in an exam or in a meeting.}
\end{frame}

\begin{frame}{What $6ND$ leaves out}
  \begin{columns}[T]
    \begin{column}{0.52\textwidth}
      \textbf{Attention's quadratic term.} Comparing every token with every
      other token uses \emph{no parameters}, so $6ND$ cannot see it at all. At a
      4k context window it adds about \textbf{10\%} on top. The work grows with
      the square of the context length, so at 128k context it dominates the bill.
      \vspace{0.5em}

      \textbf{Embeddings.} Usually left out of the parameter count. For a 70B
      model you can ignore them. For a 20M model with a 32k vocabulary they are
      a large fraction of the whole thing.
    \end{column}
    \begin{column}{0.48\textwidth}
      \begin{block}{Activation checkpointing}
        \small
        Save fewer intermediate activations and recompute them during the
        backward pass. Full recomputation adds roughly one forward pass:
        \[ C \approx 8ND \quad \text{rather than} \quad 6ND \]
        Selective checkpointing adds less compute. The exact trade-off depends
        on which operations are recomputed.
      \end{block}
      \vspace{0.3em}
      \small
      So $6ND$ is a \emph{lower bound on the useful work}. Measured wall-clock
      FLOPs are always higher.
    \end{column}
  \end{columns}
  \note{Teach them to say ``approximately'' out loud. The approximation is good
  to maybe twenty per cent for a conventional model at ordinary context
  lengths, and it breaks down badly at very long contexts, where attention is
  most of the bill.

  The size of the attention correction, if asked: relative to the linear layers
  it is about $n / (6\,d_{\text{model}})$, where $n$ is the context length and
  $d_{\text{model}}$ the model width. At $n = d_{\text{model}} = 4096$ that is
  roughly $1/6$ of a factor, i.e.\ ten per cent-ish of total compute; at
  $n = 128$k with the same width it is well over one.

  Activation checkpointing is the first appearance of the compute-for-memory
  trade that the whole of section four is about. Plant the eight-ND number now,
  reuse it in fifteen minutes.}
\end{frame}

\begin{frame}[shrink=9]{Worked: your Lab 1 run}
  Your model has \textbf{20 million} parameters and you train it on
  \textbf{1 billion} tokens, at six operations each:
  \[
    C \;=\; 6ND \;=\; 6 \times (2\times10^{7}) \times (1\times10^{9})
      \;=\; 1.2\times10^{17}\ \flops
  \]
  \vspace{0.2em}
  \begin{columns}[T]
    \begin{column}{0.5\textwidth}
      One GH200 does about $10^{15}$ of those operations per second at peak,
      and a good run turns perhaps 40\% of peak into useful work. At
      $4\times10^{14}$ per second you get through the whole budget in roughly
      \textbf{300 seconds}.
      \vspace{0.3em}

      That is five minutes of pure arithmetic.
    \end{column}
    \begin{column}{0.5\textwidth}
      \begin{block}{Measure this on Roihu}
        \small
        Because a 20M-parameter model \textbf{cannot keep an H100 busy}. The
        matmuls are too small, the data loader and the Python overhead
        dominate, and your real MFU will be a few per cent.
        \vspace{0.4em}
        Twenty minutes is a planning estimate, not a guarantee. Record the
        actual throughput and wall time; their gap from the arithmetic estimate
        is itself a useful result.
      \end{block}
    \end{column}
  \end{columns}
  \vspace{0.5em}
  \small
  At cloud rates of roughly two euro per GPU-hour, your entire pretraining run
  costs well under one euro.
  \source{GH200 peak bf16 throughput from NVIDIA datasheets; 40\% MFU is a
          typical figure for a well-tuned large run, not for a tiny one.}
  \note{This is the slide that makes 6ND feel like a tool rather than a fact.
  Do the arithmetic live; it is one multiplication.

  The five-versus-twenty-minutes point is the pedagogically valuable part and
  it is a genuine result from the dry run. Small models are bad at using big
  GPUs. Tell them to look at their own MFU in the lab and to not be surprised
  when it is two per cent.

  This also quietly explains why nobody trains tiny models on H100s in
  practice, and why the hardware market is shaped the way it is.}
\end{frame}

% The Lab 1 calculation is enough for the main lecture; frontier-cost detail is
% retained here for reference.
\iffalse
\begin{frame}[shrink=4]{Worked: what a real run costs}
  \centering
  \begin{tabular}{@{}L{2.5cm} r r r L{3.4cm}@{}}
    \toprule
    \textbf{Run} & \textbf{Params} & \textbf{Tokens} & \textbf{Compute (FLOPs)} & \textbf{Reported cost}\\
    \midrule
    Lab 1 Part B & 20M & 1B & $1.2\times10^{17}$ & under one euro\\[0.3em]
    Llama 3 8B & 8B & 15T & $7.2\times10^{23}$ & 39M GPU-h for the family\\[0.3em]
    Llama 3 405B & 405B & 15.6T & $3.8\times10^{25}$ & 30.8M H100-h, about \$62M\\[0.3em]
    DeepSeek V3 & 37B active & 14.8T & $3.3\times10^{24}$ & 2.79M H800-h, about \$5.6M\\
    \bottomrule
  \end{tabular}
  \vspace{0.5em}
  \begin{block}{What the \$5.6M figure does not include}
    It covers the \textbf{final run only}, and leaves out salaries, the data
    pipeline, every abandoned experiment, all post-training, and the cost of
    owning the cluster.
  \end{block}
  \source{DeepSeek-AI, \textit{DeepSeek-V3 Technical Report}, Dec 2024;
          Grattafiori et al., \textit{The Llama 3 Herd of Models}, 2024.
          Dollar figures use the papers' own assumed rental rate of \$2 per GPU-hour.}
  \note{The DeepSeek number was reported in the press as ``a frontier model for
  five million dollars'' and that framing is misleading in exactly the way the
  box describes. DeepSeek themselves say in the paper that the figure covers
  the official training run only.

  Two more things the figure leaves out, which did not fit in the box: the
  ablation runs that chose the recipe, and the capital cost of owning a cluster
  rather than renting one. A credible rule of thumb is that the final run is a
  minority of what a frontier model really costs.

  Make them feel the ratio: their run is about one twenty-seven-millionth of
  DeepSeek V3. That is the span this field operates over, and it is why the
  same equation has to work at both ends.

  Note the asterisk on DeepSeek: 37B is the \emph{active} parameter count,
  because 6ND counts the parameters each token actually touches. A good
  student will ask; have the answer ready.}
\end{frame}
\fi

\begin{frame}[shrink=12]{Why deployed models are often trained past the Chinchilla optimum}
  \centering
  \begin{tabular}{@{}L{3.2cm} r r r@{}}
    \toprule
    \textbf{Model} & \textbf{Params} & \textbf{Tokens} & \textbf{tokens / param}\\
    \midrule
    Chinchilla (2022) & 70B & 1.4T & 20\\
    Llama 3 8B (2024) & 8B & 15T & about 1\,900\\
    DeepSeek V3 (2024) & 37B active & 14.8T & about 400\\
    Your Lab 1 run & 20M & 1B & 50\\
    \bottomrule
  \end{tabular}
  \vspace{0.8em}
  \begin{gotcha}[Chinchilla-optimal is the wrong target for a model you will serve]
    Chinchilla minimises \textbf{training} compute for a given loss. If you are
    going to run this model billions of times, the quantity to minimise is
    \textbf{inference} compute, roughly two operations per parameter per
    generated token, which does not depend on how long you trained at all.
    \vspace{0.3em}
    So you overtrain deliberately: spend more now, to get the same loss from a
    permanently smaller model.
  \end{gotcha}
  \vspace{0.4em}
  \small
  Llama 3 8B used roughly \textbf{90 times} the compute that Chinchilla would
  have prescribed for an 8B model, and Meta said explicitly that the loss was
  still improving when they stopped.
  \source{Hoffmann et al.\ 2022; Grattafiori et al.\ 2024.}
  \note{This is the third gotcha and it is the one that most often appears in
  exams as ``is Chinchilla wrong?''. The answer is no. Chinchilla answers the
  question it asked, which is not the question a company shipping a product is
  asking.

  The asymmetry that decides it: you pay the training cost once and it scales
  with D, while you pay the inference cost on every request and it does not
  depend on D at all. Say that twice. It is the hinge of the argument and the
  opening claim of the week 3 lecture.

  The ninety-times figure is worth quoting precisely: Chinchilla-optimal for 8B
  is about 160 billion tokens, and Llama 3 used fifteen trillion.}
\end{frame}

% Detailed emergence discussion belongs with evaluation in week 4.
\iffalse
\begin{frame}{What scaling laws do not predict}
  \begin{columns}[T]
    \begin{column}{0.52\textwidth}
      A scaling law predicts \key{loss}, and nothing beyond it.
      \vspace{0.5em}

      Nobody can derive from the law how a smooth drop in cross-entropy turns
      into the ability to do multi-step arithmetic, follow a format or use a
      tool. We observe that afterwards.
      \vspace{0.5em}

      \textbf{Emergence.} Some benchmarks sit near chance up to a certain scale
      and then jump sharply. That is the ``emergent abilities'' claim.
    \end{column}
    \begin{column}{0.48\textwidth}
      \begin{block}{The deflationary counterpoint}
        \small
        Much apparent emergence is a \textbf{metric artefact}. Exact-match and
        multiple-choice accuracy are discontinuous: a model that is steadily
        getting better at a five-step problem scores zero until the last step
        flips.
        \vspace{0.4em}
        Score the same runs with a continuous metric, such as token edit
        distance or the log-likelihood of the correct answer, and the sharp
        jump often turns into a smooth curve.
      \end{block}
    \end{column}
  \end{columns}
  \vspace{0.4em}
  \small
  This is not settled. The honest position: discontinuous \emph{metrics} explain
  a lot of claimed emergence, and they do not obviously explain all of it.
  \source{Wei et al., \textit{Emergent Abilities of Large Language Models}, 2022;
          Schaeffer et al., \textit{Are Emergent Abilities a Mirage?}, 2023.}
  \note{Present both sides and do not resolve it, because it is not resolved.
  The Schaeffer argument is clean and largely correct about the examples it
  examines, and it does not license the claim that nothing interesting ever
  happens discontinuously.

  The transferable skill here is the one about metrics: before you believe a
  phase transition in a plot, check whether the metric is continuous. That
  habit is worth more than either position in the debate, and it comes back in
  week 4.}
\end{frame}
\fi

\begin{frame}{Loss is useful, but it is not the final goal}
  \begin{columns}[T]
    \begin{column}{0.54\textwidth}
      Scaling laws predict validation loss. Product and research decisions
      usually depend on something else: task accuracy, reliability, latency,
      memory use or cost per request.
      \vspace{0.5em}

      Two models with similar loss can differ because of their data mixture,
      tokenizer, context length and post-training. A small loss improvement may
      also be too expensive to justify at inference time.
    \end{column}
    \begin{column}{0.46\textwidth}
      \begin{block}{How scaling laws are used}
        \small
        Use small runs to forecast whether a larger pretraining run is healthy
        and what loss it may reach.
        \vspace{0.4em}
        Then evaluate the resulting model on the tasks and constraints that
        matter for its intended use.
        \vspace{0.4em}
        Loss is a training diagnostic and forecasting target; it does not
        replace downstream evaluation.
      \end{block}
    \end{column}
  \end{columns}
  \vfill
  \begin{takeaway}
    $C \approx 6ND$ estimates what a run costs. It does not tell us whether the
    resulting model is useful.
  \end{takeaway}
  \note{This prevents a common overinterpretation. Scaling laws are useful for
  planning training, but the dependent variable is usually validation loss.
  Task usefulness must be measured separately.

  Connect forward to week 4: downstream evaluation is where the choice of
  metric, held-out data and uncertainty becomes central.}
\end{frame}

% =====================================================================
\section{Training at scale}
% =====================================================================

\begin{frame}{Why training uses much more memory than inference}
  \centering
  \begin{tabular}{@{}L{4.2cm} C{2.6cm} C{2.6cm}@{}}
    \toprule
    \textbf{What you must hold} & \textbf{Inference} & \textbf{Training}\\
    \midrule
    parameters, bf16 & 2 bytes & 2 bytes\\
    gradients, bf16 & --- & 2 bytes\\
    master parameters, fp32 & --- & 4 bytes\\
    Adam moments $m$ and $v$, fp32 & --- & 8 bytes\\
    \midrule
    \textbf{per parameter} & \textbf{2 bytes} & \textbf{16 bytes}\\
    \bottomrule
  \end{tabular}
  \vspace{0.4em}

  \raggedright
  Training also stores \key{activations}, which depend on batch size, sequence
  length and model depth.
  \vspace{0.25em}

  \key{Useful rule of thumb: 16 bytes per parameter} for this mixed-precision
  AdamW setup. Exact requirements vary, but fitting for inference does not mean
  fitting for training.
  \source{Standard mixed-precision AdamW accounting; cf.\ Rajbhandari et al.,
          \textit{ZeRO}, 2020.}
  \note{Make them write down sixteen. It is the companion number to 6ND and
  between them they cover most back-of-envelope questions in this field.

  The fp32 master copy surprises people. The reason is that a bf16 parameter
  has about three significant decimal digits, so adding a small update to a
  large weight silently does nothing. You keep the authoritative copy in fp32
  and cast down for the matmuls.

  The last line is the sentence to repeat: the eight-times gap between serving
  and training is why LoRA exists, and that is next week.}
\end{frame}

\begin{frame}[shrink=3]{Memory examples: 7B and 20M parameters}
  \begin{columns}[T]
    \begin{column}{0.5\textwidth}
      \begin{block}{A 7B model, full fine-tune or pretrain}
        \[ 7\times10^{9} \times 16\ \text{bytes} \approx 112\ \text{GB} \]
        \vspace{-0.6em}
        \begin{itemize}\small
          \item one GH200 has \textbf{96\,GiB}, about 103\,GB
          \item so the \emph{optimiser state alone} does not fit, before a
            single activation is stored
          \item you can serve this model on half a card, and training it needs
            rather more than one
        \end{itemize}
      \end{block}
    \end{column}
    \begin{column}{0.5\textwidth}
      \begin{block}{Your Lab 1 model}
        \[ 2\times10^{7} \times 16\ \text{bytes} \approx 320\ \text{MB} \]
        \vspace{-0.6em}
        \begin{itemize}\small
          \item that is \textbf{0.3\%} of the card
          \item you can use an enormous batch, and you should
          \item your bottleneck will be the data loader and kernel launch
            overhead, not memory
        \end{itemize}
      \end{block}
    \end{column}
  \end{columns}
  \vspace{0.6em}
  \centering
  The entire rest of this section exists because of the left-hand box.
  \vfill
  \small
  Three honest escape routes from the left-hand box: train something smaller;
  shard the state across many GPUs; or do not train all the parameters. We do
  the second one today and the third one next week.
  \note{The contrast is the pedagogical point, so put both boxes up at once and
  let them sit for a moment.

  Students consistently underestimate how much harder training is than
  inference, because every tutorial they have read loads a model for inference
  and it just works. One hundred and twelve gigabytes against ninety-six is the
  cleanest way to make that concrete on hardware they will actually touch.

  The three escape routes are the roadmap for the rest of the section and for
  week 2. Say them and then start on sharding.}
\end{frame}

\begin{frame}{Activations, and buying memory with compute}
  \begin{columns}[T]
    \begin{column}{0.52\textwidth}
      The backward pass needs the intermediate tensors the forward pass
      produced, the \key{activations}. Stored naively, the total is about
      \[
        \text{batch} \times \text{sequence length} \times \text{model width}
          \times \text{layers}
      \]
      \emph{In words:} one number stored for every slot, in every layer, for
      every token in the batch.
      \vspace{0.4em}

      Activation memory grows with batch size, sequence length, model width and
      depth. It becomes especially important when the context length increases.
    \end{column}
    \begin{column}{0.48\textwidth}
      \begin{block}{Activation checkpointing}
        \small
        Store activations only at layer boundaries. Recompute the rest during
        the backward pass.
        \vspace{0.4em}
        \textbf{Memory}: save selected layer boundaries rather than every
        intermediate tensor.\\[0.3em]
        \textbf{Compute}: recompute the missing intermediates during backward.
        Full recomputation adds roughly one forward pass; selective methods add
        less.
        \vspace{0.4em}
        Selective variants recompute only what is cheap to recompute and
        expensive to store, which usually means attention rather than the
        feed-forward block.
      \end{block}
    \end{column}
  \end{columns}
  \source{Chen et al., \textit{Training Deep Nets with Sublinear Memory Cost}, 2016;
          Korthikanti et al., \textit{Reducing Activation Recomputation}, 2022.}
  \note{The thing worth insisting on: activations are the one memory term that
  does not care how big your model is. A 1B model at 128k context can be harder
  to train than a 7B model at 2k.

  The constant in the activation formula depends on how aggressively the
  framework fuses kernels, and in practice it lands somewhere between about 10
  and 30 bytes per slot per layer. The classic Chen et al.\ result is sharper than
  ``checkpoint every block'': with the optimal placement, memory drops from
  order $L$ to order $\sqrt{L}$ in the number of layers, for the same single
  extra forward pass.

  Checkpointing is the archetypal systems trade in this field: spend compute,
  which you can buy linearly, to save memory, which you cannot. That pattern
  recurs constantly, including in the KV-cache tricks from Tuesday.}
\end{frame}

\begin{frame}[shrink=9]{Four ways to split a training run}
  \begin{columns}[T]
    \begin{column}{0.5\textwidth}
      \centering
      \textbf{\small Data parallel}\\[0.3em]
      \begin{tikzpicture}[
        gpu/.style={draw=UHBlue, fill=UHBlue!18, minimum width=8mm,
                    minimum height=8mm, font=\tiny, align=center, inner sep=1pt},
        shard/.style={draw=UHGrey!60, fill=UHLight, minimum width=8mm,
                    minimum height=4mm, font=\tiny, inner sep=1pt}]
        \foreach \i/\x in {1/0, 2/1.1, 3/2.2, 4/3.3}{
          \node[gpu] at (\x,0) {whole\\model};
          \node[shard] at (\x,-0.78) {data \i};
        }
        \draw[flow] (-0.45,0.62) -- (3.75,0.62);
        \node[lbl, font=\tiny\color{UHGrey}] at (1.65,0.88) {all-reduce gradients each step};
      \end{tikzpicture}
      \vspace{1.0em}

      \textbf{\small Tensor parallel}\\[0.3em]
      \begin{tikzpicture}[
        gpu/.style={draw=UHOrange, fill=UHOrange!15, minimum width=8mm,
                    minimum height=8mm, font=\tiny, align=center, inner sep=1pt}]
        \foreach \i/\x in {1/0, 2/1.1, 3/2.2, 4/3.3}{
          \node[gpu] at (\x,0) {$\tfrac14$ of\\each $W$};
        }
        \draw[flowhi] (-0.45,0.62) -- (3.75,0.62);
        \node[lbl, font=\tiny\color{UHGrey}] at (1.65,0.88) {all-reduce \emph{inside} every layer};
        \node[lbl, font=\tiny\color{UHGrey}] at (1.65,-0.72) {one batch, one layer, four GPUs};
      \end{tikzpicture}
    \end{column}
    \begin{column}{0.5\textwidth}
      \centering
      \textbf{\small Pipeline parallel}\\[0.3em]
      \begin{tikzpicture}[
        gpu/.style={draw=UHBlue, fill=UHBlue!18, minimum width=9mm,
                    minimum height=8mm, font=\tiny, align=center, inner sep=1pt}]
        \node[gpu] (p1) at (0,0) {L\,1--8};
        \node[gpu] (p2) at (1.25,0) {L\,9--16};
        \node[gpu] (p3) at (2.5,0) {L\,17--24};
        \node[gpu] (p4) at (3.75,0) {L\,25--32};
        \draw[flow] (p1) -- (p2); \draw[flow] (p2) -- (p3); \draw[flow] (p3) -- (p4);
        \node[lbl, font=\tiny\color{UHGrey}] at (1.9,-0.72) {activations forward, gradients back};
      \end{tikzpicture}
      \vspace{1.0em}

      \textbf{\small Sequence / context parallel}\\[0.3em]
      \begin{tikzpicture}[
        seg/.style={draw=UHGreen, fill=UHGreen!12, minimum width=11mm,
                    minimum height=8mm, font=\tiny, align=center, inner sep=1pt}]
        \node[seg] at (0,0) {tok\\1--32k};
        \node[seg] at (1.25,0) {tok\\32--64k};
        \node[seg] at (2.5,0) {tok\\64--96k};
        \node[seg] at (3.75,0) {tok\\96--128k};
        \node[lbl, font=\tiny\color{UHGrey}] at (1.9,-0.78) {all-gather the keys and values};
      \end{tikzpicture}
    \end{column}
  \end{columns}
  \vspace{0.5em}
  \raggedright
  \small
  Two words will keep coming back. An \key{all-reduce} adds one tensor up
  across every GPU and hands the total back to all of them, and an
  \key{all-gather} gives every GPU a copy of every other GPU's piece. You pay
  for splitting in exactly that traffic. The four axes are \key{orthogonal},
  and a large run uses three or four at once.
  \note{Put the four up together before taking them one at a time, exactly as
  with Tuesday's KV taxonomy. The organising question for each is the same:
  \emph{what is replicated, what is split, and what has to be communicated.}

  Define the collectives out loud the first time, because this audience has no
  distributed-computing background: GPUs cannot read each other's memory, so
  every exchange is an explicit collective operation over the network. All-reduce
  = everybody contributes, everybody receives the sum. All-gather = everybody
  contributes, everybody receives the concatenation.

  Say the orthogonality point now and the composition point at the end of the
  section, so they have somewhere to put each slide as it arrives.}
\end{frame}

\begin{frame}[shrink=12]{Data parallel, and ZeRO}
  \begin{columns}[T]
    \begin{column}{0.52\textwidth}
      \textbf{Plain data parallel.} Every GPU holds the whole model and
      processes a different slice of the batch; the gradients are then summed
      across all GPUs before the step.
      \begin{itemize}\small
        \item trivially simple, almost perfectly scalable in compute
        \item \alert{but every GPU still holds all 16 bytes per parameter}
      \end{itemize}
      \vspace{0.5em}

      \textbf{ZeRO / FSDP.} To \key{shard} something is to give each GPU only
      its own slice of it and fetch the rest over the network when needed. ZeRO
      shards the two optimiser averages, then the gradients, then the weights
      themselves, collecting each layer's weights only for the moment that
      layer is in use.
    \end{column}
    \begin{column}{0.48\textwidth}
      \centering
      \begin{tabular}{@{}L{2.3cm} L{2.1cm} C{1.35cm}@{}}
        \toprule
        \textbf{Stage} & \textbf{What is split} & \textbf{GB/GPU}\\
        \midrule
        none & nothing & 112\\
        ZeRO-1 & averages & 38\\
        ZeRO-2 & $+$ gradients & 26\\
        ZeRO-3 / FSDP & $+$ weights & 14\\
        \bottomrule
      \end{tabular}

      \vspace{0.6em}
      \raggedright\small
      A 7B model across \textbf{8} GPUs. Only the last row fits on one card,
      and it pays for that by moving far more bytes over the network.
    \end{column}
  \end{columns}
  \source{Rajbhandari et al., \textit{ZeRO}, 2020; PyTorch FSDP, 2023.}
  \note{The key insight of ZeRO is embarrassingly simple and worth saying
  plainly: in plain data parallelism every GPU holds an identical copy of state
  it is not using, and that is pure waste. Sharding it changes nothing
  mathematically.

  The general formulas behind the table, with $N$ parameters over $P$ GPUs:
  unsharded is $16N$; ZeRO-1 shards the fp32 master copy and the two moment
  buffers, giving $4N + 12N/P$; ZeRO-2 also shards gradients, $2N + 14N/P$;
  ZeRO-3 shards everything, $16N/P$. The table instantiates these at
  $N = 7\times10^9$, $P = 8$.

  Communication, if asked: plain data parallel is one all-reduce of $N$ values
  per step. ZeRO-3 adds an all-gather of each layer's weights on the way forward
  and again on the way back, so roughly 50\% more traffic.

  FSDP is what they would actually use, and it is ZeRO-3 under a different
  name. Mention that at their project scale, a 4B model with LoRA, none of this
  is needed.

  If someone asks why anyone still uses ZeRO-1: because communication is not
  free and stage 1 gets most of the memory saving for almost no extra traffic.}
\end{frame}

% Detailed distributed-systems mechanisms are optional reference material.
\iffalse
\begin{frame}{Tensor parallel}
  \begin{columns}[T]
    \begin{column}{0.54\textwidth}
      Split the \emph{weight matrices} themselves across GPUs, inside each
      layer.
      \vspace{0.4em}

      In a feed-forward block, each GPU gets a slice of \emph{columns} of the
      first weight matrix and the matching slice of \emph{rows} of the second.
      Each then computes a partial output from its own slice, and one
      \key{all-reduce}, which adds those partials up across GPUs, recovers the
      true output. Attention splits the same way, a group of heads per GPU.
      \vspace{0.5em}

      \textbf{Cost}: two all-reduces per transformer block, in the critical path
      of every forward and backward pass. The GPUs cannot proceed without each
      other.
    \end{column}
    \begin{column}{0.46\textwidth}
      \begin{block}{When to use it}
        \small
        Only \textbf{within} a node, over NVLink. Tensor parallelism across a
        slower interconnect is a reliable way to make a cluster slower than a
        single machine.
        \vspace{0.4em}
        Typical configuration: tensor parallel degree 8 inside a node, then
        pipeline or data parallel between nodes.
        \vspace{0.4em}
        It is also the main reason \texttt{num\_attention\_heads} is usually a
        multiple of 8.
      \end{block}
    \end{column}
  \end{columns}
  \source{Shoeybi et al., \textit{Megatron-LM}, 2019.}
  \note{The column-then-row split is the one piece of actual maths in this
  section and it is worth drawing on the board. The block computes
  $Y = \mathrm{SiLU}(X W_1) W_2$. Write $W_1 = [W_1^{(1)} \mid W_1^{(2)}]$, a
  split by columns, and $W_2$ split correspondingly by rows. Then GPU $k$ holds
  $W_1^{(k)}$ and $W_2^{(k)}$ and computes
  $\mathrm{SiLU}(X W_1^{(k)}) W_2^{(k)}$; the true $Y$ is the sum of those
  partials, which is exactly one all-reduce.

  Why that order: splitting the first matrix by columns means each GPU owns
  whole output units, so the nonlinearity can be applied independently. Split
  the other way and you need a communication step \emph{before} the activation
  function, which doubles the number of collectives.

  The practical rule is the sentence in the box: tensor parallelism stays inside
  a node. They will meet this again in week 3 as a serving decision, where the
  same rule applies for the same reason.}
\end{frame}

\begin{frame}[shrink=9]{Pipeline parallel, and the bubble}
  \centering
  \begin{tikzpicture}[cell/.style={tokcell, minimum size=4.6mm}]
    \foreach \d in {1,...,4}{
      \foreach \t in {1,...,14}{
        \node[cell, fill=UHGrey!12] at (\t*0.52, -\d*0.55) {};
      }
      \node[lbl, font=\tiny\color{UHGrey}, anchor=east] at (0.3, -\d*0.55) {GPU \d};
    }
    \foreach \d in {1,...,4}{
      \foreach \m in {1,...,4}{
        \pgfmathtruncatemacro{\tf}{\m + \d - 1}
        \node[cell, fill=UHBlue!35] at (\tf*0.52, -\d*0.55) {F\m};
        \pgfmathtruncatemacro{\tb}{11 + \m - \d}
        \node[cell, fill=UHOrange!35] at (\tb*0.52, -\d*0.55) {B\m};
      }
    }
    \node[lbl, font=\tiny\color{UHGrey}] at (3.9, -2.75) {time $\longrightarrow$};
  \end{tikzpicture}
  \vspace{0.6em}

  \begin{columns}[T]
    \begin{column}{0.5\textwidth}
      \small
      Four stages, four micro-batches. Each GPU is busy for 8 of the 14 time
      slots; the grey cells are the \key{bubble}, idle hardware waiting for a
      stage that is not ready yet.
      \[
        \text{idle fraction} = \frac{p-1}{m+p-1} = \frac{3}{7} \approx 43\%
      \]
      \emph{In words:} with $p$ stages and $m$ micro-batches, you lose the time
      it takes to fill the pipe and drain it again.
    \end{column}
    \begin{column}{0.5\textwidth}
      \small
      \textbf{The fix is more micro-batches.} With 32 micro-batches over 4
      stages the bubble falls to about 9\%; cleverer schedules shrink it further.
      \vspace{0.4em}

      \textbf{The attraction}: the only thing sent between GPUs is the
      activations at stage boundaries, so pipeline parallelism survives a slow
      network link.
    \end{column}
  \end{columns}
  \source{Huang et al., \textit{GPipe}, 2019; Narayanan et al., \textit{PipeDream} / Megatron-LM, 2021.}
  \note{Walk the diagram left to right with a finger. GPU 1 does all four
  forward micro-batches, then has nothing to do until gradients come back from
  the far end of the pipe.

  The formula is worth memorising because it tells you the design rule
  immediately: the bubble shrinks when micro-batches outnumber stages, so deep
  pipelines need large batches. Interleaved and 1F1B schedules do better still
  and also reduce peak activation memory, by starting a backward pass as soon as
  one micro-batch is ready rather than waiting for all of them.

  Then the contrast that organises the section: tensor parallel needs a fast
  interconnect, pipeline parallel tolerates a slow one, and that is most of how
  people choose between them.}
\end{frame}

\begin{frame}[shrink=13]{Sequence parallel, and 3D parallelism}
  \begin{columns}[T]
    \begin{column}{0.5\textwidth}
      \textbf{Sequence / context parallel.} Split the \emph{tokens} of one
      sequence across GPUs. Each device owns a contiguous chunk and its
      activations, which is the only way to train at 128k context without the
      activation memory exploding.
      \vspace{0.4em}

      Attention is the hard part, because every token must still be able to look
      at every earlier token. Ring Attention passes blocks of keys and values
      around the ring of GPUs so that the communication overlaps with the
      computation and costs almost nothing.
      \vspace{0.5em}

      \textbf{Sequence-parallel norms} are the cheap cousin: the norm layers can
      be split along the sequence too, removing duplicated activation memory for
      free.
    \end{column}
    \begin{column}{0.5\textwidth}
      \begin{block}{How a real run composes them}
        \small
        A typical frontier configuration:
        \vspace{0.3em}
        \textbf{tensor} parallel 8 --- within a node, over NVLink\\[0.2em]
        \textbf{pipeline} parallel 8--16 --- across nodes\\[0.2em]
        \textbf{data} parallel / FSDP --- across everything else\\[0.2em]
        \textbf{context} parallel --- added only for the long-context phase\\[0.2em]
        \textbf{expert} parallel --- if the model is an MoE
      \end{block}
      \vspace{0.3em}
      \small
      Llama 3 405B used 4D parallelism across up to 16\,384 H100s.
    \end{column}
  \end{columns}
  \source{Liu et al., \textit{Ring Attention}, 2023;
          Grattafiori et al., \textit{The Llama 3 Herd of Models}, 2024.}
  \note{Context parallelism is the newest of the four and the one most tied to
  the long-context arms race. Note that Llama 3 turned it on only for the
  long-context stage, which is a nice concrete illustration that the
  parallelism configuration changes during a run.

  The composition box is what to photograph. The ordering is not arbitrary: the
  most communication-hungry axis goes on the fastest link, and so on outward.
  From that one principle they can reconstruct the configuration themselves.}
\end{frame}
\fi

\begin{frame}[shrink=16]{Precision}
  A number format splits its bits between the \key{exponent}, which sets how
  large or small a value it can represent, and the \key{mantissa}, which sets
  how many significant digits it keeps.

  \vspace{0.5em}
  \centering
  \begin{tabular}{@{}l c c L{5.0cm}@{}}
    \toprule
    \textbf{Format} & \textbf{Exponent} & \textbf{Mantissa} & \textbf{Why it matters}\\
    \midrule
    fp32 & 8 & 23 & master weights and optimiser state\\[0.3em]
    fp16 & 5 & 10 & \alert{overflows}; needs loss scaling\\[0.3em]
    bf16 & 8 & 7 & \good{fp32's range}, fewer digits\\[0.3em]
    fp8 (E4M3) & 4 & 3 & 2$\times$ throughput; per-block scaling\\
    \bottomrule
  \end{tabular}

  \vspace{0.6em}
  \raggedright
  \key{Why bf16 became common:} it keeps fp32's exponent range, which reduces
  overflow and underflow problems compared with fp16, while retaining
  half-precision storage and matrix operations.
  \vspace{0.3em}

  DeepSeek V3 used \key{fp8} for selected operations, with higher precision
  where needed. The potential speed-up depends on hardware, kernels and the
  scaling strategy.
  \source{DeepSeek-AI, \textit{DeepSeek-V3 Technical Report}, Dec 2024;
          Micikevicius et al., \textit{Mixed Precision Training}, 2018.}
  \note{The exponent-versus-mantissa framing is the whole slide. Say it as a
  slogan: in deep learning you need dynamic range, not decimal places.

  Two things the table had to drop. fp32 is also the reference format that
  everything else is measured against, and bf16 needs no loss scaling at all,
  which is most of why people switched to it. fp8 is still not routine outside
  the largest labs.

  Loss scaling is the standard fp16 workaround: multiply the loss before the
  backward pass so small gradients remain representable, then divide the
  gradients before the update. bf16 reduces the need for this because it has a
  wider exponent range, although implementations still need numerical checks.

  On fp8: the interesting part is how much care it takes. DeepSeek used
  fine-grained per-tile scaling with fp32 accumulation at the sensitive
  operations. Their report is unusually explicit about which operations they
  kept in higher precision, and it is good reading.}
\end{frame}

\begin{frame}{MFU: how much of the hardware is doing model work}
  \key{Model FLOPs utilisation (MFU)} compares estimated model arithmetic with
  the peak arithmetic the allocated GPUs could have performed:
  \[
    \mathrm{MFU} \;=\;
      \frac{6ND}{\text{run time}\;\times\;\text{number of GPUs}
                 \;\times\;\text{peak rate}}
  \]
  \emph{In words:} how much of the available compute became useful model
  operations?
  \vspace{0.5em}
  \begin{columns}[T]
    \begin{column}{0.5\textwidth}
      \textbf{Why it is below 100\%}
      \begin{itemize}\small
        \item GPUs wait for data and communication
        \item norms and softmax do not reach matrix-multiply peak
        \item checkpointing repeats work
        \item small matrices do not fill the GPU
      \end{itemize}
    \end{column}
    \begin{column}{0.5\textwidth}
      \begin{block}{What it is used for}
        \small
        Track MFU while tuning one training system. A sudden drop can reveal a
        data-loader, communication or kernel problem.
        \vspace{0.4em}
        Comparisons across models, hardware or numeric formats need care because
        both $6ND$ and ``peak rate'' are approximations.
      \end{block}
    \end{column}
  \end{columns}
  \note{The point is not a universal ``good'' percentage. MFU is most useful as
  a within-system diagnostic: hold the model and hardware fixed, make one
  change, and check whether useful throughput improves. Tiny Lab 1 models will
  have low utilisation because their matrix operations are too small.}
\end{frame}

\begin{frame}{Large training runs fail regularly}
  \begin{columns}[T]
    \begin{column}{0.52\textwidth}
      Llama 3 405B, over a \textbf{54-day} pretraining window on 16k H100s:
      \begin{itemize}\small
        \item \textbf{466} job interruptions
        \item \textbf{419} of them unexpected
        \item \textbf{78\%} of the unexpected ones were hardware --- GPU
          failures most of all
        \item effective training time still above \textbf{90\%}, because the
          recovery was automated
        \item power draw across the datacentre swinging by tens of megawatts as
          thousands of GPUs ramped together
      \end{itemize}
    \end{column}
    \begin{column}{0.48\textwidth}
      \begin{block}{What this looks like from inside}
        \small
        Silent data corruption that shows up only as a loss that is slightly
        too high.\\[0.3em]
        NaN hunts that take three days and end at one bad node.\\[0.3em]
        Checkpoint cadence as an explicit trade: a 405B optimiser state is
        about 6.5\,TB, so writing it costs real minutes.\\[0.3em]
        On LUMI, the same class of problems plus an immature software stack
        that absorbed a large share of the project's engineering.
      \end{block}
    \end{column}
  \end{columns}
  \source{Grattafiori et al., \textit{The Llama 3 Herd of Models}, 2024;
          personal experience, Poro / Viking on LUMI, 2023--2024.}
  \note{This slide tells them what the job is actually like. Tell it warmly
  rather than reading the list. Most of what happens on a large pretraining run
  is operations work.

  The number to quote is 419 unexpected interruptions in 54 days, roughly eight
  a day, on the best-engineered run that has been documented publicly.

  The LUMI anecdote is first-hand, so take thirty seconds. What surprised me
  most was how much of the calendar went into making the software stack work at
  all on hardware nobody else was training LLMs on.

  Clock check: should be around 17:15 here.}
\end{frame}

% This repeats the compute/cost comparison and is kept only as reference.
\iffalse
\begin{frame}[shrink=14]{What a run costs, in three currencies}
  \centering
  \begin{tabular}{@{}L{2.6cm} L{2.6cm} L{2.3cm} L{3.2cm}@{}}
    \toprule
    \textbf{Run} & \textbf{GPU-hours} & \textbf{Wall clock} & \textbf{Money}\\
    \midrule
    Llama 3 405B & 30.8M H100-h & 54 days on 16k GPUs & about \$62M at \$2/h\\[0.3em]
    DeepSeek V3 & 2.79M H800-h & about 2 months on 2k GPUs & about \$5.6M at \$2/h\\[0.3em]
    Poro 34B & a few hundred LUMI GPUs & several weeks & a national allocation\\[0.3em]
    Your Lab 1 run & about 0.3 GPU-h & 20 minutes & under one euro\\
    \bottomrule
  \end{tabular}

  \vspace{0.9em}
  \raggedright
  \small
  The scarcest of the three is usually \key{wall clock}.
  You cannot buy a fifty-four-day run in a week by adding GPUs indefinitely,
  because communication costs grow and the bubble grows with them. That is why
  the architecture gets frozen months before anybody sees the model.
  \vspace{0.4em}

  Meta reported about 11\,390 tonnes of CO$_2$-equivalent for training the
  Llama 3 family, on a location-based accounting.
  \source{Grattafiori et al.\ 2024; DeepSeek-AI 2024. Dollar figures use each
          paper's own assumed rental rate and are not what the owners paid.}
  \note{The wall-clock point is the one practitioners feel first. Scaling to
  more GPUs has sharply diminishing returns, so a big model takes months no
  matter how rich you are, and that fixes a long lead time into every decision.
  It is exactly the Poro retrospective from earlier.

  Mention the carbon figure without editorialising. It is in the paper, it is
  the only number of its kind published at that scale, and the honest framing is
  that inference emissions across a model's lifetime dwarf it.

  Then move on briskly; mid-training is short and section six needs its ten
  minutes for the demo.}
\end{frame}
\fi

% =====================================================================
\section{Mid-training}
% =====================================================================

\begin{frame}{Mid-training: targeted continued pretraining}
  \centering
  \begin{tikzpicture}[node distance=4mm]
    \node[mod, minimum width=24mm] (pre) {pretraining\\{\tiny trillions of tokens, general web}};
    \node[modhi, minimum width=24mm, right=of pre] (mid) {mid-training\\{\tiny $10^{10}$--$10^{11}$ tokens, specialised}};
    \node[mod, minimum width=24mm, right=of mid] (post) {post-training\\{\tiny SFT, DPO, RLVR}};
    \draw[flow] (pre) -- (mid);
    \draw[flow] (mid) -- (post);
  \end{tikzpicture}
  \vspace{1.0em}

  \begin{columns}[T]
    \begin{column}{0.5\textwidth}
      \textbf{Same objective, different data.} Mid-training is still
      next-token prediction with AdamW. What changes is that the mixture is
      chosen deliberately rather than scraped.
    \end{column}
    \begin{column}{0.5\textwidth}
      \textbf{Why it exists.} The cheapest place to install a capability is
      usually the last few per cent of pretraining, rather than the start of a
      trillion-token run or a fine-tune on a hundred thousand examples.
    \end{column}
  \end{columns}
  \vspace{0.6em}
  \centering
  \small
  The term only stabilised around 2024--2025. Before that the same work was
  called ``annealing'', ``the decay phase'', ``continued pretraining'' or nothing at all.
  \source{OLMo 2 (Allen AI, 2024) is the clearest publicly documented example;
          Raschka, \textit{The State of LLMs}, 2025, on the terminology.}
  \note{Flag the terminology explicitly, because students reading papers from
  2022 to 2026 will meet four different names for the same stage and assume
  they are four different things.

  The definition to give is operational: if the loss is cross-entropy and the
  optimiser is AdamW and the data is curated rather than crawled, it is
  mid-training. Everything else is a judgement call.

  Say up front that this section is three ideas --- long context, annealing,
  capability injection --- so they know it is short.}
\end{frame}

\begin{frame}[shrink=8]{Long-context extension: the problem}
  Pretraining at 128k context is wasteful: attention is quadratic, activation
  memory is linear in sequence length, and almost no documents are that long.
  So you pretrain short, at 4k or 8k, and extend afterwards.
  \vspace{0.6em}
  \begin{columns}[T]
    \begin{column}{0.52\textwidth}
      \textbf{Why you cannot just feed it more tokens.} RoPE encodes position by
      rotating each attention dimension by an angle proportional to the token's
      position. Think of a bank of clocks, each ticking at a different speed,
      with the whole range set by a number called the \key{base}, usually
      $10^4$.
      \vspace{0.4em}

      Past the positions seen in training, the fast clocks have wrapped around
      many times into a combination of readings the model has never encountered.
      Attention degenerates.
    \end{column}
    \begin{column}{0.48\textwidth}
      \begin{block}{Base scaling, the blunt fix}
        \small
        Increase the base: $10^4 \to 5\times10^5$ (Llama 3) or higher. Every
        frequency slows down, so the same position index maps to a smaller
        angle and more positions fit inside the range the model knows.
        \vspace{0.4em}
        It works, but you pay for it in resolution on \emph{local} position,
        which is what the model leans on most.
      \end{block}
    \end{column}
  \end{columns}
  \source{Su et al., RoFormer, 2021; Grattafiori et al.\ 2024.}
  \note{The wrap-around intuition is the one to give: RoPE is a set of clocks
  running at different speeds, and extrapolating past training length means
  reading a clock configuration that never occurred during training.

  The formula behind the clock metaphor, if anyone wants it: dimension pair $i$
  of a head of width $d$ rotates at angular frequency
  $\theta_i = \text{base}^{-2i/d}$, so $i = 0$ is the fastest clock and
  $i = d/2$ the slowest. Raising the base stretches the whole bank towards
  slower rotation.

  Base scaling is the one-line fix and it is what most models actually ship.
  Make sure they see what it costs. Slowing all the clocks down also slows the
  fast ones that encode ``the previous few tokens'', and that is where most of
  the useful signal sits.}
\end{frame}

% Position-interpolation details are deferred to the inference lecture.
\iffalse
\begin{frame}{YaRN, and the two routes to long context}
  \begin{columns}[T]
    \begin{column}{0.5\textwidth}
      \key{YaRN} applies base scaling selectively:
      \begin{itemize}\small
        \item leave the \textbf{high-frequency} dimensions alone, so local
          position stays sharp
        \item interpolate the \textbf{low-frequency} dimensions, which carry
          long-range position
        \item nudge the sharpness of attention, because spreading it over a much
          longer context otherwise flattens it
        \item then fine-tune briefly on long documents
      \end{itemize}
      \vspace{0.4em}
      Typically buys a 4--8$\times$ extension for a small fraction of a per cent
      of the pretraining compute.
    \end{column}
    \begin{column}{0.5\textwidth}
      \begin{block}{Route A: extend after the fact}
        \small
        \textbf{Qwen3 235B-A22B}: 32k native, \textbf{131k with YaRN}.
      \end{block}
      \begin{block}{Route B: build it into the architecture}
        \small
        \textbf{Qwen3-Next 80B-A3B}: \textbf{262k native}, which is affordable
        only because three of every four attention layers are linear-attention
        blocks with a context-independent state size.
      \end{block}
      \vspace{0.3em}
      \small
      You are looking at the same decision Tuesday's attention slides
      described, taken at a different point in the pipeline.
    \end{column}
  \end{columns}
  \source{Peng et al., \textit{YaRN}, 2023; Raschka,
          \textit{A Dream of Spring for Open-Weight LLMs}, Feb 2026.}
  \note{The two-route framing is what makes this slide worth its place. A
  config that says 262k native and a config that says 32k-extended-to-131k
  describe completely different engineering, and the difference is visible in
  the attention design from Tuesday.

  Practical consequence they will hit: a YaRN-extended model is usually
  configured with the scaling factor in config.json, and applying it when you
  do not need long context makes short-context quality slightly worse. Several
  stacks therefore enable it only on request.}
\end{frame}
\fi

\begin{frame}[shrink=8]{Annealing on a higher-quality mixture}
  \begin{columns}[T]
    \begin{column}{0.54\textwidth}
      During the decay phase of a WSD schedule, swap the data mixture for a
      small, very high-quality one and let the learning rate fall to near zero.
      \vspace{0.5em}

      \textbf{Why use the decay phase?} It provides a short, inexpensive stage
      in which the data mixture can be changed and compared. A useful intuition
      is that smaller late updates refine the model without repeating the full
      pretraining run; the exact mechanism is still an empirical question.
      \vspace{0.5em}

      \textbf{OLMo 2} is the documented example: a curated mixture of
      high-quality web, maths and instruction-like data, applied over the final
      50B to 300B tokens, then several annealed runs averaged into one model.
    \end{column}
    \begin{column}{0.46\textwidth}
      \begin{block}{Why WSD and annealing fit together}
        \small
        A cosine schedule has no flat phase to branch from, so annealing means
        re-planning the whole run.
        \vspace{0.4em}
        With WSD you keep one stable-phase checkpoint and anneal it several
        times over, varying the mixture, the budget or the target language, for
        a few per cent of the cost each.
        \vspace{0.4em}
        This flexibility is an important reason WSD became common.
      \end{block}
    \end{column}
  \end{columns}
  \source{OLMo 2 (Allen AI, 2024); Hu et al., \textit{MiniCPM}, 2024.}
  \note{Call back to the schedule plot from section one and point at the final
  decay. This is where that corner of the curve earns its keep.

  The mechanism is intuitive and roughly right: what the model sees while the
  learning rate is small has outsized influence. The same mechanism is behind
  catastrophic forgetting in week 2, so plant that connection here.

  The reusable-checkpoint argument is the one that persuades practitioners. You
  pay for one stable phase and get several finished models out of it.}
\end{frame}

\begin{frame}{What targeted data can add}
  What actually goes into a mid-training mixture, and why:

  \vspace{0.4em}
  \centering
  \begin{tabular}{@{}L{3.6cm} L{7.2cm}@{}}
    \toprule
    \textbf{Ingredient} & \textbf{What it is for}\\
    \midrule
    maths & symbolic step chains; the best lever on reasoning\\[0.25em]
    code & structure and verifiable chains; lifts reasoning\\[0.25em]
    long documents & teaches the model to \emph{use} its context window\\[0.25em]
    instruction-like text & question-answer formatting for SFT to grip\\[0.25em]
    target language & the cheapest place to add a new language\\[0.25em]
    synthetic textbooks & dense explanatory prose; narrows the distribution\\
    \bottomrule
  \end{tabular}

  \vspace{0.5em}
  \raggedright
  \small
  By the bottom two rows, mid-training is hard to tell apart from SFT without
  the loss masking.
  \source{OLMo 2 (2024); Aryabumi et al., \textit{To Code or Not To Code?}, 2024;
          Gunasekar et al., \textit{Textbooks Are All You Need}, 2023.}
  \note{Two rows deserve a gloss beyond what is on the slide. Code lifts
  general reasoning, not only coding ability, which is why labs put it in even
  when they do not care about code. Synthetic textbooks work well and they also
  narrow the model's output distribution, so a model trained on too many of
  them writes in one register.

  The row worth dwelling on is instruction-like text. Putting
  question-answer formatted documents into the \emph{pretraining} mixture makes
  the later SFT stage dramatically more effective, because the format is
  already in the model's distribution and SFT only has to select it rather than
  install it.

  That observation is also the cleanest preview of next Tuesday's central
  claim: post-training selects behaviour that pretraining made possible.}
\end{frame}

% Terminology detail retained for reference; the main lecture uses the
% operational definition on the opening mid-training slide.
\iffalse
\begin{frame}[shrink=10]{Where mid-training ends and post-training begins}
  \begin{columns}[T]
    \begin{column}{0.5\textwidth}
      \textbf{Where people draw the line}
      \begin{itemize}\small
        \item loss on all tokens $\Rightarrow$ mid-training;
          loss on assistant tokens only $\Rightarrow$ SFT
        \item documents $\Rightarrow$ mid-training;
          templated conversations $\Rightarrow$ SFT
        \item learning rate still decaying from the pretraining schedule
          $\Rightarrow$ mid-training
      \end{itemize}
      \vspace{0.5em}
      \textbf{Why the line will not hold.} All three criteria are conventions,
      and some published recipe violates each of them. Nothing in the
      mathematics tells the stages apart.
    \end{column}
    \begin{column}{0.5\textwidth}
      \begin{block}{What to take from this}
        \small
        When you read a model report, do not ask which stage something is in.
        Ask three concrete questions:
        \vspace{0.3em}
        \emph{what data?} \quad \emph{what loss, on which tokens?} \quad
        \emph{what learning rate?}
        \vspace{0.4em}
        Those three answers determine the model. The stage names are
        bookkeeping.
      \end{block}
    \end{column}
  \end{columns}
  \vfill
  \begin{takeaway}
    Mid-training is where a general model becomes \emph{your} model, at one per
    cent of the cost of pretraining. That is why so much of the field's
    attention has moved here since 2024.
  \end{takeaway}
  \note{The three questions in the box are the transferable part and they are
  worth writing on the board. They work for every training stage in this
  course, including everything next Tuesday.

  Resist the temptation to give a crisp definition of mid-training. There is
  not one, and pretending otherwise would teach them to argue about vocabulary
  instead of reading recipes.

  Clock check: should be around 17:27 here. The demo needs ten minutes and it
  is the best part of the lecture, so protect it.}
\end{frame}
\fi

% =====================================================================
\section{What pretraining buys}
% =====================================================================

\begin{frame}[fragile]{In-context learning, live}
  \demo{a base model, with no instruction tuning and no chat template.}
\begin{lstlisting}
Review: The food was cold and the waiter ignored us.  Sentiment: negative
Review: Beautiful room, spotless, and the staff were kind.  Sentiment: positive
Review: It arrived late and the box was crushed.  Sentiment: negative
Review: Exactly what I wanted, and cheaper than the shop.  Sentiment:
\end{lstlisting}
  \vspace{0.3em}
  The model completes \texttt{\ positive}. Nobody trained it to classify
  sentiment, and no weight was updated.
  \vspace{0.5em}

  It picked up sentiment classification \textbf{from the prompt}, because the
  prompt is a document whose most likely continuation is the correct label.
  \vspace{0.4em}

  \key{In-context learning} is the emergent capability that made the whole
  prompting era possible. You get it free from next-token prediction on a
  corpus that happens to contain lists, tables and worked examples.
  \source{Brown et al., \textit{Language Models are Few-Shot Learners}, 2020.}
  \note{Run this live on Roihu. The machine is up today, maintenance is not
  until 3 and 4 November. Use a genuine base model, not an instruct model, and
  say so, because the whole point collapses if someone thinks this is chat.

  The sentence that lands: the model is continuing a document, and in that
  document the next thing after the fourth ``Sentiment:'' is a label.
  Classification falls out of continuation.

  If the few-shot prompt fails on the first try, that is also a useful lesson.
  Base models are sensitive to formatting, and a stray newline changes the
  answer.}
\end{frame}

\begin{frame}{What in-context learning means}
  \begin{columns}[T]
    \begin{column}{0.5\textwidth}
      \textbf{What happens.} The prefix changes the model's \emph{conditional
      distribution}. No weights move, and the effect disappears as soon as the
      context does.
      \vspace{0.5em}

      \textbf{What to be careful about.} Calling it learning is a stretch.
      There is a real literature analysing it as gradient descent inside the
      forward pass, but the case is suggestive rather than settled.
      \vspace{0.5em}

      \textbf{Empirically}: performance rises with the number of examples, with
      the consistency of the \emph{format}, and surprisingly weakly with whether
      the labels in the examples are actually correct.
    \end{column}
    \begin{column}{0.5\textwidth}
      \begin{block}{Why it emerges at all}
        \small
        The pretraining corpus is full of documents in which a pattern is
        established and then continued: glossaries, tables, translation pairs,
        exercise-and-answer sets, code with consistent conventions.
        \vspace{0.4em}
        A model that models those documents well \emph{must} learn to infer a
        rule from a prefix and apply it. There was never a separate objective
        for this.
      \end{block}
      \vspace{0.3em}
      \small
      In Lab 1 Part D you run the same few-shot prompt on your tiny model and on
      a larger base model. Treat this as a qualitative comparison: scale,
      architecture, tokenizer and training data all differ.
    \end{column}
  \end{columns}
  \source{Brown et al.\ 2020; Min et al., \textit{Rethinking the Role of
          Demonstrations}, 2022.}
  \note{The label-correctness result is the fun one: shuffling the labels in
  few-shot demonstrations degrades performance much less than you would expect,
  which suggests the demonstrations are mostly specifying the \emph{task and
  the format} rather than teaching the mapping.

  Use that to set up the honest position: we have good empirical
  characterisations of in-context learning and no settled mechanistic account.
  This is a live research area and a reasonable thesis topic.

  Point at Part D. Their 20M model will show essentially no in-context learning,
  and that negative result is the whole purpose of the exercise.}
\end{frame}

\begin{frame}[fragile]{Now ask it a question}
  \demo{the same base model, same session, no few-shot examples.}
\begin{lstlisting}
Prompt:  What is the capital of Finland?

Output:  What is the capital of Sweden?
         What is the capital of Norway?
         What is the largest city in Finland?
         What is the population of Helsinki?
         Answer the following questions in complete sentences.
         1. What is the capital of Finland? ____________
         2. What is ...
\end{lstlisting}
  \vspace{0.3em}
  Nothing has necessarily gone wrong here. More questions are one plausible
  continuation because question lists and worksheets occur in the pretraining
  corpus.
  \vspace{0.4em}

  \alert{And it will not stop.} No token in the corpus reliably means ``the
  assistant has finished its turn'', because there is no assistant.
  \note{This is the pivot of the whole week and it deserves a pause. Let the
  output run on the screen for a few seconds longer than is comfortable.

  The framing that matters: what we trained and what we wanted are two
  different things. The model is a very good model of documents. We wanted an
  assistant, and we never asked for one.

  Then name the two specific missing things, turn-taking and stopping, because
  those are mechanical, concrete, and fixed in exactly one way, which is
  Tuesday's lecture.}
\end{frame}

\begin{frame}[shrink=10]{What the base model does not have}
  \centering
  \begin{tabular}{@{}L{4.2cm} L{6.6cm}@{}}
    \toprule
    \textbf{Missing} & \textbf{Why pretraining could not supply it}\\
    \midrule
    stopping & no end-of-turn token exists in web documents\\[0.25em]
    turn-taking & no notion of ``user'' and ``assistant'' roles\\[0.25em]
    instruction following & instructions occur, but obeying them is rarely the
      likeliest continuation\\[0.25em]
    a preferred answer & good and bad answers appear at their web
      frequencies\\[0.25em]
    refusal & declining a request is rare, and inconsistent where it
      occurs\\[0.25em]
    a consistent persona & the corpus is millions of different authors\\
    \bottomrule
  \end{tabular}

  \vspace{0.6em}
  \raggedright
  \small
  Every row is a problem of \textbf{selection} rather than \textbf{knowledge}.
  The few-shot demo showed the base model can produce the right behaviour; it
  just does not reach for it by default. That is the thesis of next Tuesday.
  \source{Ouyang et al., \textit{InstructGPT}, 2022, for the original framing.}
  \note{Say the selection-versus-knowledge sentence slowly. It is the opening
  claim of lecture three and the thing students most often get backwards. They
  arrive believing fine-tuning teaches the model facts, when what it mostly
  does is select and shape behaviour pretraining already made possible.

  On the fourth row, spell out that web frequency is a poor proxy for quality:
  the corpus is full of confident wrong answers, and a base model will happily
  reproduce them at the rate they appear.

  The evidence for the claim is already on the previous two slides: the same
  weights that rambled at a direct question classified sentiment correctly when
  the format made the behaviour likely. The capability was there. The default
  was wrong.}
\end{frame}

% Detailed metric discussion is deferred to the evaluation lecture.
\iffalse
\begin{frame}{Perplexity against downstream accuracy}
  \begin{columns}[T]
    \begin{column}{0.52\textwidth}
      Within one model family, on one tokenizer, over one data mixture, lower
      validation loss reliably means better downstream performance. That is why
      loss is a usable training signal at all.
      \vspace{0.5em}

      \textbf{They come apart when:}
      \begin{itemize}\small
        \item the tokenizers differ, so the unit itself changes
        \item the data mixtures differ, and some corpora are simply easier
        \item the benchmark is near chance at that scale, so it measures noise
        \item the capability is bottlenecked by \emph{format} rather than by
          knowledge
        \item the evaluation corpus is in the training set
      \end{itemize}
    \end{column}
    \begin{column}{0.48\textwidth}
      \begin{block}{The practical rule}
        \small
        Use loss to decide whether a \textbf{run is healthy}.
        \vspace{0.4em}
        Use downstream evaluation to decide whether a \textbf{model is good}.
        \vspace{0.4em}
        Do not substitute one for the other, and do not compare losses across
        tokenizers. That was the first gotcha of the lecture, ninety minutes
        ago.
      \end{block}
      \vspace{0.3em}
      \small
      Week 4 is entirely about the second half of that rule.
    \end{column}
  \end{columns}
  \note{Close the loop with the first gotcha explicitly. Students respond well
  to a callback and it signals that the lecture was built rather than
  assembled.

  The format-bottleneck bullet is the one that will bite them in week 2: a
  model that knows the answer and will not emit it in the expected format
  scores zero, and its loss will look fine.

  Keep this brief. One more slide in this section.}
\end{frame}
\fi

\begin{frame}[shrink=11]{What pretraining gives us}
  \begin{columns}[T]
    \begin{column}{0.52\textwidth}
      \textbf{What pretraining produced}
      \begin{itemize}\small
        \item a compressed model of a large slice of human text
        \item world knowledge, syntax, code semantics, arithmetic, discourse
          structure
        \item the ability to infer a task from a prefix and continue it
        \item a representation good enough that a few thousand examples can
          redirect its behaviour
      \end{itemize}
      \vspace{0.5em}
      \textbf{What it did not produce}
      \begin{itemize}\small
        \item reliable assistant-style turn-taking, instruction following or
          stopping behaviour
      \end{itemize}
    \end{column}
    \begin{column}{0.48\textwidth}
      \begin{block}{Why base models are still released}
        \small
        A base model is a useful starting point when researchers want to study
        or shape behaviour without an existing instruction-tuning policy.
        \vspace{0.4em}
        Lab 2 compares a base model with a corresponding instruct model. You
        may use a different suitable pair from the one used elsewhere in your
        project.
      \end{block}
    \end{column}
  \end{columns}
  \vfill
  \begin{takeaway}
    Pretraining supplies broad language and task capability. Post-training can
    select, strengthen and sometimes add behaviours using much less compute,
    but it still depends on what the base model can represent.
  \end{takeaway}
  \note{Use the takeaway to distinguish broad capability learned from large
  corpora from the more targeted behaviour shaping in post-training. The
  distinction is useful, but not absolute: later training can add knowledge and
  capabilities too.

  Remind them that this is exactly why the course spends one lecture on
  pretraining, which they will probably never do, and two on post-training and
  inference, which they certainly will.

  Clock check: should be around 17:37 here. Lab briefing next, eight minutes,
  then out on time.}
\end{frame}

% =====================================================================
\section{Lab 1 briefing}
% =====================================================================

\begin{frame}[shrink=12]{Lab 1: from tokenizer to tiny GPT}
  \begin{columns}[T]
    \begin{column}{0.56\textwidth}
      \textbf{A. Tokenizer} \hfill \textit{CPU Jupyter, about 1 h}\\
      Byte-level BPE on a Finnish Wikipedia slice at two or three vocabulary
      sizes, then compare fertility against Qwen3.5 and Gemma 4 on FLORES-200.
      \vspace{0.3em}

      \textbf{B. Tiny GPT} \hfill \textit{GPU batch, about 20 min / run}\\
      Train a 10--30M-parameter decoder with \emph{your} tokenizer. Script and
      Slurm template provided.
      \vspace{0.3em}

      \textbf{C. Scaling} \hfill \textit{3 GPU batch jobs}\\
      Three model sizes at a fixed token budget. Plot validation loss against
      compute, estimating compute as $6ND$.
      \vspace{0.3em}

      \textbf{D. Tiny versus pretrained} \hfill \textit{GPU Jupyter, about 20 min}\\
      The same few-shot prompt on your tiny model and on a real base model.
      \vspace{0.3em}

      \textbf{E. Project} \hfill \textit{planning}\\
      Choose a task, model and held-out set. You may revise this later.
    \end{column}
    \begin{column}{0.44\textwidth}
      \begin{block}{How the parts connect to the lectures}
        \small
        the training loop and the loss curve $\Rightarrow$ B\\[0.25em]
        $C \approx 6ND$ and scaling laws $\Rightarrow$ C\\[0.25em]
        in-context learning $\Rightarrow$ D\\[0.25em]
        and Tuesday's tokenization $\Rightarrow$ A
      \end{block}
      \vspace{0.3em}
      \small
      A complete run of the three supplied model sizes should use roughly
      \textbf{one GPU-hour} in total.
      \vspace{0.5em}

      Start with \texttt{lab1-example/README.md}; it includes a CPU smoke test,
      the training script, configurations and plotting code.
    \end{column}
  \end{columns}
  \note{In Part A, run the fertility comparison in English, Finnish and one
  more language of their choice. Part B's provided script is about 300 lines,
  so they should read it rather than treat it as a black box.

  Map each part back to a slide from this week out loud. It makes the lab feel
  like a consequence of the lectures rather than a separate chore.

  Recommend doing A before B, then C and D after one training run works. Part E
  can be drafted at any point and revised later.

  One GPU-hour each is worth saying because every year somebody is afraid of
  using the allocation and runs nothing.}
\end{frame}

% The runnable template is in lab1-example/slurm/train.slurm.
\iffalse
\begin{frame}[fragile]{The Slurm template}
\begin{lstlisting}[language=slurm]
#!/bin/bash
#SBATCH --account=project_2021893
#SBATCH --partition=gputest         # debug here first; 15 min cap
#SBATCH --nodes=1
#SBATCH --gres=gpu:1                # one GH200
#SBATCH --cpus-per-task=16
#SBATCH --mem=64G
#SBATCH --time=00:15:00
#SBATCH --output=logs/%x-%j.out

module purge
module load pytorch

srun python train_tiny_gpt.py \
     --config configs/medium.yaml --compile
\end{lstlisting}
  \vspace{0.2em}
  \small
  For Part C, use \texttt{small.yaml}, \texttt{medium.yaml} and
  \texttt{large.yaml}. Keep the tokenizer, token budget and seed fixed. After
  the smoke test, use \texttt{gpumedium} and request enough time. \quad
  \texttt{sbatch} \textbullet\ \texttt{squeue -{}-me} \textbullet\
  \texttt{seff <jobid>}
  \note{The rule from Tuesday, repeated because it is the one that saves them
  the most time: debug on gputest, with a tiny token budget, until the script
  runs end to end. Only then move to gpumedium.

  Point at the seed flag. Fixed seeds are the difference between a scaling plot
  that means something and three runs that differ for unknown reasons.

  Point out what changes in Part C: layer count, width and head count. The data,
  tokenizer, token budget and random seed stay fixed. That makes this a
  controlled scaling comparison rather than three unrelated runs.}
\end{frame}
\fi

\begin{frame}{How to check that a run is working}
  \centering
  \begin{tabular}{@{}L{3.4cm} L{3.4cm} L{4.0cm}@{}}
    \toprule
    \textbf{Check} & \textbf{Expected pattern} & \textbf{If it differs}\\
    \midrule
    first-step loss & close to $\ln(V)$; about 10.4 for 32k tokens &
      check initialization, vocabulary and labels\\[0.25em]
    training loss & falls quickly, then more slowly & inspect the learning rate
      and batches\\[0.25em]
    validation loss & follows training loss & check the split and possible
      overfitting\\[0.25em]
    throughput & becomes stable after warmup & inspect the data loader and GPU
      utilisation\\[0.25em]
    output files & metrics, summary and weights exist & read the end of the job
      log\\
    \bottomrule
  \end{tabular}

  \vspace{0.6em}
  \raggedright
  \small
  \key{The one-line sanity check}: your initial loss should be close to the log
  of your vocabulary size, and $\ln 32000 \approx 10.4$. If your run starts
  far from that value, investigate before spending GPU time.
  \note{A randomly initialised model often predicts close to uniformly, making
  $\ln(V)$ a useful approximate sanity check. It is not a proof that target
  alignment is correct: a shifted target can produce a similar initial loss.
  Stronger checks are that loss stays finite, falls on a short run and can
  overfit a tiny repeated batch. The reference implementation tests the full
  pipeline and falling finite loss on a bundled fixture.}
\end{frame}

\begin{frame}{Material for the final report}
  \begin{columns}[T]
    \begin{column}{0.5\textwidth}
      \begin{enumerate}
        \item a \key{fertility table} covering your tokenizers and two
          production tokenizers, across three languages
        \item \key{loss curves} for all your runs, on one set of axes
        \item a \key{scaling plot} of validation loss against $6ND$, log--log
        \item one paragraph comparing \key{your run's compute} with a published
          frontier run, using $6ND$ for both
        \item your project plan: task, model and held-out set
      \end{enumerate}
    \end{column}
    \begin{column}{0.5\textwidth}
      \begin{block}{What makes the analysis useful}
        \small
        A number with the method that produced it, and a sentence saying what
        you concluded from it.
        \vspace{0.4em}
        If a result surprises you, investigate it and record what you checked.
        \vspace{0.4em}
        ``The three runs did not order the way I expected, and here is what I
        checked'' is a good paragraph.
      \end{block}
      \vspace{0.3em}
      \small
      Keep the scripts, configurations and machine-readable result files in
      your repository. There is no separate Lab 1 submission.
    \end{column}
  \end{columns}
  \note{Item four is the one that connects the lecture to the lab most
  directly, and it is the one most often skipped. They should produce two FLOP
  figures with the same equation and then say something honest about the ratio.

  Repeat the grading principle from Tuesday: we grade measured results and the
  decisions drawn from them, rather than effort described. The lab is small
  enough that everyone can get the numbers, so what they say about them is what
  separates the reports.}
\end{frame}

\begin{frame}{Working on Lab 1}
  \begin{columns}[T]
    \begin{column}{0.52\textwidth}
      \begin{itemize}
        \item \alert{Roihu is down for maintenance on 3--4 November.} Those are
          not available for GPU work.
        \item \key{There is no supervised lab this week.} The next time we are
          all in a room with GPUs is 5 November.
        \item Questions go to the course channel. It is watched all week and
          answers are visible to everybody.
        \item There is no intermediate submission. Keep your work for the final
          report and repository.
      \end{itemize}
    \end{column}
    \begin{column}{0.48\textwidth}
      \begin{block}{If something breaks}
        \small
        Post the \textbf{job id}, the \textbf{command}, and the \textbf{last
        twenty lines} of the log. Three quarters of problems are solved from
        that alone.
        \vspace{0.4em}
        Do not post screenshots of a terminal.
      \end{block}
      \vspace{0.3em}
      \begin{block}{A useful starting point}
        Run the bundled CPU smoke test, then Part A. This checks your
        environment and data flow before you request GPU time.
      \end{block}
    \end{column}
  \end{columns}
  \note{Mention the maintenance window and post it in the course channel.
  Students can do the tokenizer work and inspect the implementation without a
  GPU.

  The bug-report format is worth insisting on from week one. It trains a habit
  that makes the next six weeks much cheaper for everybody, and it is the same
  information a colleague would want.}
\end{frame}

% =====================================================================
% A second title-only summary slide is not needed in a 90-minute lecture.
\iffalse
\begin{frame}[standout]
  Three things to remember
  \note{Say these three cleanly and slowly. They are reused verbatim in the
  week 7 wrap-up, and they are the examinable summary of the lecture.

  Then one deep breath before the hand-off slide, which is thirty seconds and
  sets up Tuesday.}
\end{frame}
\fi

\begin{frame}{Three things to remember}
  \begin{enumerate}
    \item \key{$C \approx 6ND$ gives a useful first estimate of training
      compute.} Two FLOPs per parameter per token forward and about four
      backward. Remember what the approximation leaves out.
      \vspace{0.6em}
    \item \key{Data choices strongly affect model quality.} Inspect the data
      for obvious problems, then compare pipeline choices by training small,
      controlled models and evaluating them.
      \vspace{0.6em}
    \item \key{A deployed model may be trained past the compute-optimal point.}
      Extra training can produce a smaller model with lower inference cost,
      trading a one-time expense for repeated savings.
  \end{enumerate}
  \vfill
  \small
  And one number to carry: \textbf{16 bytes per parameter} to train in mixed
  precision, against 2 to serve.
  \note{Deliver these as three sentences, not as a slide to be read. Pause
  after each.

  If anyone asks which one matters most, say the second. The first is
  arithmetic and the third is a business decision, while the second is a
  research method they can apply to any dataset question for the rest of their
  careers.}
\end{frame}

\begin{frame}{Tuesday: the model that will not stop talking}
  \begin{columns}[T]
    \begin{column}{0.54\textwidth}
      We ended on a base model that answered a question with more questions and
      then kept going.
      \vspace{0.5em}

      Nobody patches that. It is the gap between a model of \emph{documents}
      and an \emph{assistant}, and closing it takes a whole discipline:
      \vspace{0.3em}

      \centering
      SFT \textrightarrow{} preference optimisation \textrightarrow{}
      RLVR \textrightarrow{} distillation

      \vspace{0.3em}
      \raggedright
      Each stage exists because the previous one has a specific limitation.
    \end{column}
    \begin{column}{0.46\textwidth}
      \begin{block}{Tuesday 3 November}
        \small
        Post-training. Chat templates, loss masking, DPO, GRPO, and why being
        able to \emph{check} an answer removed the need for a reward model.

        \vspace{0.4em}
        The base-versus-instruct demo will be pre-recorded, because Roihu is
        down that day.
      \end{block}
      \vspace{0.3em}
      \small
      \textbf{This week}: begin Lab 1 and keep the outputs for your final
      report. There is no intermediate submission.
    \end{column}
  \end{columns}
  \note{Thirty seconds. End on the image of the rambling base model, because it
  is the thing they will remember and it is the question Tuesday answers.

  Remind them once more about the maintenance window, and then let them go on
  time.}
\end{frame}

\appendix

\begin{frame}[shrink=18]{Reading}
  \begin{columns}[T]
    \begin{column}{0.5\textwidth}
      \textbf{Core}
      \begin{itemize}\small
        \item Stanford CS336, \textit{Language Modeling from Scratch}, the
          lectures on data and on systems
        \item Hugging Face, \textit{The Smol Training Playbook}, 2025
        \item Karpathy, \textit{nanoGPT}
      \end{itemize}
      \vspace{0.4em}
      \textbf{Data}
      \begin{itemize}\small
        \item Penedo et al., \textit{The FineWeb Datasets}, 2024
        \item Lozhkov et al., \textit{FineWeb-Edu}, 2024
        \item Soldaini et al., \textit{Dolma}, 2024
        \item Luukkonen et al., \textit{Poro 34B and the Blessing of
          Multilinguality}, 2024
      \end{itemize}
    \end{column}
    \begin{column}{0.5\textwidth}
      \textbf{Scaling}
      \begin{itemize}\small
        \item Kaplan et al.\ 2020; Hoffmann et al.\ 2022 (Chinchilla)
        \item Muennighoff et al., \textit{Scaling Data-Constrained Language
          Models}, 2023
        \item Schaeffer et al., \textit{Are Emergent Abilities a Mirage?}, 2023
      \end{itemize}
      \vspace{0.4em}
      \textbf{Systems}
      \begin{itemize}\small
        \item Rajbhandari et al., \textit{ZeRO}, 2020
        \item Shoeybi et al., \textit{Megatron-LM}, 2019
        \item Grattafiori et al., \textit{The Llama 3 Herd of Models}, 2024;
          read section 3 on infrastructure
        \item DeepSeek-AI, \textit{DeepSeek-V3 Technical Report}, 2024, for fp8
          training in detail
      \end{itemize}
    \end{column}
  \end{columns}
  \vspace{0.3em}
  \small
  For current developments, Raschka's \textit{Ahead of AI}:
  \textit{A Dream of Spring for Open-Weight LLMs} (Feb 2026) and
  \textit{The State of LLMs} (2025).
  \note{If they read one thing, make it section 3 of the Llama 3 paper. It is
  the most honest published account of what large-scale training actually
  involves, and it is readable.

  The Smol Training Playbook is the best practical companion to this lecture
  for anyone who wants to train something real at a scale a university can
  afford.}
\end{frame}


\end{document}
